% Cloudra 源程序清单（软著申请用），tectonic 直接编译
\documentclass[10pt,a4paper]{article}
\usepackage{ctex}
\usepackage[top=16mm,bottom=16mm,left=20mm,right=20mm]{geometry}
\usepackage{listings}
\usepackage{fancyhdr}
\usepackage{fontspec}
\setmonofont{DejaVu Sans Mono}
\setCJKmonofont{WenQuanYi Zen Hei}
\lstset{
  basicstyle=\ttfamily\fontsize{6.5}{8.3}\selectfont,
  columns=fullflexible,
  keepspaces=true,
  breaklines=true,
  breakatwhitespace=false,
  breakindent=0pt,
  postbreak=,
  tabsize=4,
  frame=none,
  xleftmargin=0pt,
  xrightmargin=0pt,
  aboveskip=0pt,
  belowskip=0pt,
}
\pagestyle{fancy}
\fancyhf{}
\fancyhead[L]{\small\ttfamily Cloudra V2.0 源程序}
\fancyhead[R]{\small\ttfamily 第 \thepage{} 页}
\renewcommand{\headrulewidth}{0.4pt}
\setlength{\headheight}{14pt}
\begin{document}
\setcounter{page}{1}
\begin{lstlisting}[firstnumber=1,numbers=left,numberstyle=\tiny\ttfamily,numbersep=5pt]
cloudra/__init__.py:1	"""Cloudra - AI Agent collaboration framework for chemical research."""
cloudra/__init__.py:2	
cloudra/__init__.py:3	from importlib import import_module
cloudra/__init__.py:4	
cloudra/__init__.py:5	__version__ = "2.0.0-beta"
cloudra/__init__.py:6	
cloudra/__init__.py:7	_LAZY_EXPORTS = {
cloudra/__init__.py:8	    "Agent": ("cloudra.core.agent", "Agent"),
cloudra/__init__.py:9	    "Coordinator": ("cloudra.core.agent", "Coordinator"),
cloudra/__init__.py:10	}
cloudra/__init__.py:11	# 兼容垫片：包结构重排（core/services/persistence/base/infra）前的
cloudra/__init__.py:12	# `from cloudra import X` 老路径全部继续可用（惰性加载，不提前拉起重依赖）。
cloudra/__init__.py:13	_LAZY_MODULES = {
cloudra/__init__.py:14	    # core
cloudra/__init__.py:15	    "agent": "cloudra.core.agent",
cloudra/__init__.py:16	    "supervisor": "cloudra.core.supervisor",
cloudra/__init__.py:17	    # services
cloudra/__init__.py:18	    "session_store": "cloudra.services.session_store",
cloudra/__init__.py:19	    "nodes": "cloudra.services.nodes",
cloudra/__init__.py:20	    # persistence
cloudra/__init__.py:21	    "database": "cloudra.persistence.database",
cloudra/__init__.py:22	    "artifacts": "cloudra.persistence.artifacts",
cloudra/__init__.py:23	    "memory": "cloudra.persistence.memory",
cloudra/__init__.py:24	    "message_history": "cloudra.core.message_history",
cloudra/__init__.py:25	    "paths": "cloudra.persistence.paths",
cloudra/__init__.py:26	    "fs_utils": "cloudra.persistence.fs_utils",
cloudra/__init__.py:27	    "workspace": "cloudra.persistence.workspace",
cloudra/__init__.py:28	    # base
cloudra/__init__.py:29	    "session_class": "cloudra.base.session_class",
cloudra/__init__.py:30	    "skill_class": "cloudra.base.skill_class",
cloudra/__init__.py:31	    "tools_class": "cloudra.base.tools_class",
cloudra/__init__.py:32	    # infra
cloudra/__init__.py:33	    "config": "cloudra.infra.config",
cloudra/__init__.py:34	    "llm_client": "cloudra.infra.llm_client",
cloudra/__init__.py:35	    "streaming": "cloudra.infra.streaming",
cloudra/__init__.py:36	    "token_estimator": "cloudra.infra.token_estimator",
cloudra/__init__.py:37	    "web_search_native": "cloudra.infra.web_search_native",
cloudra/__init__.py:38	    "slurm": "cloudra.infra.slurm",
cloudra/__init__.py:39	    "compute_utils": "cloudra.infra.compute_utils",
cloudra/__init__.py:40	    "operations": "cloudra.infra.operations",
cloudra/__init__.py:41	    "yolo_guard": "cloudra.infra.yolo_guard",
cloudra/__init__.py:42	    "platform_utils": "cloudra.infra.platform_utils",
cloudra/__init__.py:43	    "runtime_profile": "cloudra.infra.runtime_profile",
cloudra/__init__.py:44	}
cloudra/__init__.py:45	
cloudra/__init__.py:46	__all__ = [*_LAZY_EXPORTS, *_LAZY_MODULES, "__version__"]
cloudra/__init__.py:47	
cloudra/__init__.py:48	
cloudra/__init__.py:49	def __getattr__(name: str):
cloudra/__init__.py:50	    """Keep legacy package exports without importing the Agent runtime eagerly."""
cloudra/__init__.py:51	    target = _LAZY_EXPORTS.get(name)
cloudra/__init__.py:52	    if target is not None:
\end{lstlisting}
\clearpage
\begin{lstlisting}[firstnumber=53,numbers=left,numberstyle=\tiny\ttfamily,numbersep=5pt]
cloudra/__init__.py:53	        module_name, attribute = target
cloudra/__init__.py:54	        value = getattr(import_module(module_name), attribute)
cloudra/__init__.py:55	        globals()[name] = value
cloudra/__init__.py:56	        return value
cloudra/__init__.py:57	    module_path = _LAZY_MODULES.get(name)
cloudra/__init__.py:58	    if module_path is not None:
cloudra/__init__.py:59	        value = import_module(module_path)
cloudra/__init__.py:60	        globals()[name] = value
cloudra/__init__.py:61	        return value
cloudra/__init__.py:62	    raise AttributeError(f"module 'cloudra' has no attribute {name!r}")
cloudra/__init__.py:63	
cloudra/__init__.py:64	
cloudra/__init__.py:65	def __dir__() -> list[str]:
cloudra/__init__.py:66	    return sorted(set(globals()) | set(__all__))
cloudra/__main__.py:1	"""Allow running as `python -m cloudra`."""
cloudra/__main__.py:2	from cloudra.cli import main
cloudra/__main__.py:3	
cloudra/__main__.py:4	if __name__ == "__main__":
cloudra/__main__.py:5	    main()
cloudra/base/__init__.py:1	"""cloudra.base 子包。"""
cloudra/base/session_class.py:1	
cloudra/base/session_class.py:2	import os
cloudra/base/session_class.py:3	from datetime import datetime
cloudra/base/session_class.py:4	
cloudra/base/session_class.py:5	from cloudra.persistence.fs_utils import atomic_write_json
cloudra/base/session_class.py:6	
cloudra/base/session_class.py:7	
cloudra/base/session_class.py:8	class Session:
cloudra/base/session_class.py:9	    """内存态会话镜像；数据库（session 表）是唯一持久化事实来源。
cloudra/base/session_class.py:10	
cloudra/base/session_class.py:11	    默认不再写 per-session JSON 快照文件；调试时可设置环境变量
cloudra/base/session_class.py:12	    ``CLOUDRA_DEBUG_SESSION_JSON=1`` 恢复写文件行为（upload_json 仍可显式调用）。
cloudra/base/session_class.py:13	    """
cloudra/base/session_class.py:14	    id: str
cloudra/base/session_class.py:15	    path: str
cloudra/base/session_class.py:16	    create_time: datetime
cloudra/base/session_class.py:17	    state: str = 'sleeping'
cloudra/base/session_class.py:18	
cloudra/base/session_class.py:19	    def __init__(self, id: str, path: str) -> None:
cloudra/base/session_class.py:20	        self.create_time = datetime.now()
cloudra/base/session_class.py:21	        self.id = id
cloudra/base/session_class.py:22	        self.path = path
cloudra/base/session_class.py:23	        self.state = 'sleeping'
cloudra/base/session_class.py:24	        if os.environ.get("CLOUDRA_DEBUG_SESSION_JSON") == "1":
cloudra/base/session_class.py:25	            dir_name = os.path.dirname(self.path)
cloudra/base/session_class.py:26	            if dir_name:
cloudra/base/session_class.py:27	                os.makedirs(dir_name, exist_ok=True)
cloudra/base/session_class.py:28	            self.upload_json()
cloudra/base/session_class.py:29	
cloudra/base/session_class.py:30	
cloudra/base/session_class.py:31	    def to_json(self) -> dict:
cloudra/base/session_class.py:32	        create_time = self.create_time
\end{lstlisting}
\clearpage
\begin{lstlisting}[firstnumber=33,numbers=left,numberstyle=\tiny\ttfamily,numbersep=5pt]
cloudra/base/session_class.py:33	        if isinstance(create_time, datetime):
cloudra/base/session_class.py:34	            create_time = create_time.isoformat()
cloudra/base/session_class.py:35	        json_dict = {
cloudra/base/session_class.py:36	            'id': self.id,
cloudra/base/session_class.py:37	            'path': self.path,
cloudra/base/session_class.py:38	            'create_time': create_time,
cloudra/base/session_class.py:39	            'state': self.state,
cloudra/base/session_class.py:40	        }
cloudra/base/session_class.py:41	        return json_dict
cloudra/base/session_class.py:42	
cloudra/base/session_class.py:43	    def upload_json(self):
cloudra/base/session_class.py:44	        atomic_write_json(self.path, self.to_json())
cloudra/base/session_class.py:45	
cloudra/base/skill_class.py:1	import hashlib
cloudra/base/skill_class.py:2	import os
cloudra/base/skill_class.py:3	
cloudra/base/skill_class.py:4	import frontmatter as frt
cloudra/base/skill_class.py:5	
cloudra/base/skill_class.py:6	# 模块级缓存按 SKILL.md 内容 revision 失效。Phase 6 后 skill manifest
cloudra/base/skill_class.py:7	# 允许在进程运行期间更新，不能再假设目录在整个进程生命周期中静态不变。
cloudra/base/skill_class.py:8	_SKILL_DATAS_CACHE: dict[str, tuple[str, list[dict]]] = {}
cloudra/base/skill_class.py:9	
cloudra/base/skill_class.py:10	# NOTE: Skills 模型仅作为 skill frontmatter 的字段参照，当前未被实例化使用。
cloudra/base/skill_class.py:11	# 若后续需要校验 skill 元数据，可取消注释并补全默认值。
cloudra/base/skill_class.py:12	# class Skills(BaseModel):
cloudra/base/skill_class.py:13	#     name: str
cloudra/base/skill_class.py:14	#     description: str
cloudra/base/skill_class.py:15	#     path: str
cloudra/base/skill_class.py:16	#     scripts: list = []
cloudra/base/skill_class.py:17	#     references: list = []
cloudra/base/skill_class.py:18	
cloudra/base/skill_class.py:19	
cloudra/base/skill_class.py:20	class skill_loader:
cloudra/base/skill_class.py:21	
cloudra/base/skill_class.py:22	    # Skills are bundled with the package, not user data（base/ 上溯一级到包根）
cloudra/base/skill_class.py:23	    base_path: str = os.path.join(
cloudra/base/skill_class.py:24	        os.path.dirname(os.path.dirname(os.path.abspath(__file__))), 'skills'
cloudra/base/skill_class.py:25	    )
cloudra/base/skill_class.py:26	
cloudra/base/skill_class.py:27	    def __init__(self):
cloudra/base/skill_class.py:28	        # FIX: 原代码直接用了模块级变量 base_path（缺少 self.），在实例方法中访问类属性应加 self
cloudra/base/skill_class.py:29	        self.namelst = os.listdir(self.base_path)
cloudra/base/skill_class.py:30	
cloudra/base/skill_class.py:31	        self.pathlst = [os.path.join(self.base_path, skill_name) for skill_name in self.namelst]
cloudra/base/skill_class.py:32	
cloudra/base/skill_class.py:33	        self.md_pathlst = [os.path.join(path,'SKILL.md') for path in self.pathlst if os.path.isdir(path)]
cloudra/base/skill_class.py:34	
cloudra/base/skill_class.py:35	    def skill_datas_get(self) -> list[dict]:
cloudra/base/skill_class.py:36	        revision_hash = hashlib.sha256()
cloudra/base/skill_class.py:37	        for md in sorted(self.md_pathlst):
cloudra/base/skill_class.py:38	            revision_hash.update(md.encode("utf-8"))
cloudra/base/skill_class.py:39	            try:
\end{lstlisting}
\clearpage
\begin{lstlisting}[firstnumber=40,numbers=left,numberstyle=\tiny\ttfamily,numbersep=5pt]
cloudra/base/skill_class.py:40	                with open(md, "rb") as stream:
cloudra/base/skill_class.py:41	                    revision_hash.update(stream.read())
cloudra/base/skill_class.py:42	            except OSError:
cloudra/base/skill_class.py:43	                continue
cloudra/base/skill_class.py:44	        revision = revision_hash.hexdigest()
cloudra/base/skill_class.py:45	        cached = _SKILL_DATAS_CACHE.get(self.base_path)
cloudra/base/skill_class.py:46	        if cached is None or cached[0] != revision:
cloudra/base/skill_class.py:47	            mds = self.md_pathlst
cloudra/base/skill_class.py:48	            datas = []
cloudra/base/skill_class.py:49	            for md in mds:
cloudra/base/skill_class.py:50	                with open(md, 'r', encoding='utf-8') as f:
cloudra/base/skill_class.py:51	                    post = frt.load(f)
cloudra/base/skill_class.py:52	                    metadata = post.metadata
cloudra/base/skill_class.py:53	                    content = post.content
cloudra/base/skill_class.py:54	
cloudra/base/skill_class.py:55	                    # FIX: dict.update() 返回 None，原代码把 None 当数据 append 了
cloudra/base/skill_class.py:56	                    # 这里直接修改 metadata dict 后 append（frontmatter 的 metadata 是普通 dict）
cloudra/base/skill_class.py:57	                    metadata['content'] = content
cloudra/base/skill_class.py:58	                    datas.append(metadata)
cloudra/base/skill_class.py:59	
cloudra/base/skill_class.py:60	            _SKILL_DATAS_CACHE[self.base_path] = (revision, datas)
cloudra/base/skill_class.py:61	
cloudra/base/skill_class.py:62	        return _SKILL_DATAS_CACHE[self.base_path][1]
cloudra/base/skill_class.py:63	
cloudra/base/skill_class.py:64	    def get_reference(self, skill_name, ref_name) -> str:
cloudra/base/skill_class.py:65	        # 支持多种可能的资料文件夹命名约定
cloudra/base/skill_class.py:66	        possible_dirs = ['references', 'reference', 'ref']
cloudra/base/skill_class.py:67	
cloudra/base/skill_class.py:68	        for dir_name in possible_dirs:
cloudra/base/skill_class.py:69	            ref_path = os.path.join(self.base_path, skill_name, dir_name, ref_name)
cloudra/base/skill_class.py:70	            if os.path.isfile(ref_path):
cloudra/base/skill_class.py:71	                with open(ref_path, 'r', encoding='utf-8') as f:
cloudra/base/skill_class.py:72	                    return f.read()
cloudra/base/skill_class.py:73	            elif os.path.isdir(ref_path):
cloudra/base/skill_class.py:74	                dir_filelst = os.listdir(ref_path)
cloudra/base/skill_class.py:75	                return fr'该reference为文件夹，文件夹下有：{dir_filelst} ,请根据需要输入 {ref_name}/filename 调用资料，win系统请用 {ref_name}\filename'
cloudra/base/skill_class.py:76	
cloudra/base/skill_class.py:77	        return '该资料不存在。'
cloudra/base/skill_class.py:78	
cloudra/base/skill_class.py:79	    # FIX: 原拼写为 skill_dsatas_get（多打了个 s）
cloudra/base/skill_class.py:80	    def get_content(self, skill_name) -> str:
cloudra/base/skill_class.py:81	        datas = self.skill_datas_get()
cloudra/base/skill_class.py:82	        for data in datas:
cloudra/base/skill_class.py:83	            if data['name'] == skill_name:
cloudra/base/skill_class.py:84	                return data['content']
cloudra/base/skill_class.py:85	        return 'skill不存在。'
cloudra/base/skill_class.py:86	
cloudra/base/skill_class.py:87	
cloudra/base/skill_class.py:88	
cloudra/base/skill_class.py:89	
cloudra/base/skill_class.py:90	
cloudra/base/skill_class.py:91	
\end{lstlisting}
\clearpage
\begin{lstlisting}[firstnumber=92,numbers=left,numberstyle=\tiny\ttfamily,numbersep=5pt]
cloudra/base/skill_class.py:92	
cloudra/base/skill_class.py:93	
cloudra/base/skill_class.py:94	
cloudra/base/skill_class.py:95	
cloudra/base/skill_class.py:96	
cloudra/base/skill_class.py:97	
cloudra/base/skill_class.py:98	
cloudra/base/skill_class.py:99	
cloudra/base/tools_class.py:1	from contextlib import contextmanager
cloudra/base/tools_class.py:2	from contextvars import ContextVar
cloudra/base/tools_class.py:3	from typing import TYPE_CHECKING, Any, Awaitable, Callable, Optional
cloudra/base/tools_class.py:4	
cloudra/base/tools_class.py:5	from pydantic import BaseModel
cloudra/base/tools_class.py:6	from pydantic_core import PydanticUndefined
cloudra/base/tools_class.py:7	
cloudra/base/tools_class.py:8	from cloudra.base.session_class import Session
cloudra/base/tools_class.py:9	from cloudra.persistence.database import ConversationDB, get_database_path
cloudra/base/tools_class.py:10	
cloudra/base/tools_class.py:11	# HACK: 这里原本缺少 import，属于历史遗留问题；get_tools_lst 定义在 tools/__init__.py
cloudra/base/tools_class.py:12	from cloudra.tools import get_tools_lst
cloudra/base/tools_class.py:13	from cloudra.tools.core.permission_model import Permission
cloudra/base/tools_class.py:14	
cloudra/base/tools_class.py:15	if TYPE_CHECKING:
cloudra/base/tools_class.py:16	    from cloudra.infra.yolo_guard import YoloGuard
cloudra/base/tools_class.py:17	
cloudra/base/tools_class.py:18	# runtime_v2 的审批钩子：LegacyToolExecutor 在执行工具前总是设置；
cloudra/base/tools_class.py:19	# 未设置即生产外路径（v1/CLI/Flet 已删除），request_approval 直接抛错。
cloudra/base/tools_class.py:20	# 签名: async (**fields) -> tuple[bool, str]，字段同 request_approval 的关键字参数。
cloudra/base/tools_class.py:21	_APPROVAL_HANDLER: ContextVar[
cloudra/base/tools_class.py:22	    Optional[Callable[..., Awaitable[tuple[bool, str]]]]
cloudra/base/tools_class.py:23	] = ContextVar("cloudra_approval_handler", default=None)
cloudra/base/tools_class.py:24	
cloudra/base/tools_class.py:25	
cloudra/base/tools_class.py:26	@contextmanager
cloudra/base/tools_class.py:27	def use_approval_handler(handler: Callable[..., Awaitable[tuple[bool, str]]]):
cloudra/base/tools_class.py:28	    """Scope a v2 approval handler to one tool execution (per-task isolation)."""
cloudra/base/tools_class.py:29	    token = _APPROVAL_HANDLER.set(handler)
cloudra/base/tools_class.py:30	    try:
cloudra/base/tools_class.py:31	        yield
cloudra/base/tools_class.py:32	    finally:
cloudra/base/tools_class.py:33	        _APPROVAL_HANDLER.reset(token)
cloudra/base/tools_class.py:34	
cloudra/base/tools_class.py:35	
cloudra/base/tools_class.py:36	def _tool_db_path(subagent) -> str:
cloudra/base/tools_class.py:37	    """根据 tool 关联的 Agent 返回对应 conversation 的数据库路径。"""
cloudra/base/tools_class.py:38	    conv_id = getattr(subagent, 'conversation_id', 'default') if subagent else 'default'
cloudra/base/tools_class.py:39	    return get_database_path(conv_id)
cloudra/base/tools_class.py:40	
cloudra/base/tools_class.py:41	
cloudra/base/tools_class.py:42	def get_tools() -> list[dict]:
cloudra/base/tools_class.py:43	    """将工具类加载为 OpenAI function 格式的 dict 列表。"""
cloudra/base/tools_class.py:44	    entries = get_tools_lst()
\end{lstlisting}
\clearpage
\begin{lstlisting}[firstnumber=45,numbers=left,numberstyle=\tiny\ttfamily,numbersep=5pt]
cloudra/base/tools_class.py:45	    tools_list = []
cloudra/base/tools_class.py:46	    for entry in entries:
cloudra/base/tools_class.py:47	        tool_instance = entry.instantiate()
cloudra/base/tools_class.py:48	        tools_list.append(tool_instance.tool_schema_json())
cloudra/base/tools_class.py:49	    return tools_list
cloudra/base/tools_class.py:50	
cloudra/base/tools_class.py:51	
cloudra/base/tools_class.py:52	class input_schema(BaseModel):
cloudra/base/tools_class.py:53	    '''
cloudra/base/tools_class.py:54	    工具输入模式基类，定义工具需要的输入参数和格式。
cloudra/base/tools_class.py:55	
cloudra/base/tools_class.py:56	    子类用法示例：
cloudra/base/tools_class.py:57	        class MyInput(input_schema):
cloudra/base/tools_class.py:58	            query: str = Field(description="搜索关键词")
cloudra/base/tools_class.py:59	            max_results: int = Field(default=10, description="最大返回条数")
cloudra/base/tools_class.py:60	
cloudra/base/tools_class.py:61	    通过 to_parameters_schema() 可自动生成符合 OpenAI Function Calling 规范的 parameters 字典。
cloudra/base/tools_class.py:62	    '''
cloudra/base/tools_class.py:63	
cloudra/base/tools_class.py:64	    @classmethod
cloudra/base/tools_class.py:65	    def to_parameters_schema(cls) -> dict:
cloudra/base/tools_class.py:66	        '''
cloudra/base/tools_class.py:67	        生成符合 OpenAI Function Calling 的 parameters 字典。
cloudra/base/tools_class.py:68	        基于 Pydantic model_json_schema() 微调：
cloudra/base/tools_class.py:69	        - properties / required 已由 Pydantic 根据字段声明自动生成
cloudra/base/tools_class.py:70	        - 强制附加 additionalProperties: false，防止 LLM 传入未定义字段
cloudra/base/tools_class.py:71	        '''
cloudra/base/tools_class.py:72	        schema = cls.model_json_schema()
cloudra/base/tools_class.py:73	        parameters = {
cloudra/base/tools_class.py:74	            "type": "object",
cloudra/base/tools_class.py:75	            "properties": schema.get("properties", {}),
cloudra/base/tools_class.py:76	            "required": schema.get("required", []),
cloudra/base/tools_class.py:77	            "additionalProperties": False,
cloudra/base/tools_class.py:78	        }
cloudra/base/tools_class.py:79	        if "$defs" in schema:
cloudra/base/tools_class.py:80	            parameters["$defs"] = schema["$defs"]
cloudra/base/tools_class.py:81	        return parameters
cloudra/base/tools_class.py:82	
cloudra/base/tools_class.py:83	    @classmethod
cloudra/base/tools_class.py:84	    def from_args(cls, arg_dict: dict) -> "input_schema":
cloudra/base/tools_class.py:85	        '''从 LLM 返回的 arguments dict 构造输入实例，带 Pydantic 类型校验。'''
cloudra/base/tools_class.py:86	        return cls(**arg_dict)
cloudra/base/tools_class.py:87	
cloudra/base/tools_class.py:88	    def toClass(self, arg_dict: dict) -> "input_schema":
cloudra/base/tools_class.py:89	        '''
cloudra/base/tools_class.py:90	        兼容旧接口，供 reAct.py 调用。
cloudra/base/tools_class.py:91	        实际调用链通常为 tool.input_schema.toClass(tool_args)，此时 self 为 input_schema 子类本身。
cloudra/base/tools_class.py:92	        '''
cloudra/base/tools_class.py:93	        # 当通过类调用时（如 TextReader_input_schema.toClass(args)），self 即为类对象
cloudra/base/tools_class.py:94	        if isinstance(self, type) and issubclass(self, input_schema):
cloudra/base/tools_class.py:95	            return self(**arg_dict)
cloudra/base/tools_class.py:96	        return self.__class__(**arg_dict)
\end{lstlisting}
\clearpage
\begin{lstlisting}[firstnumber=97,numbers=left,numberstyle=\tiny\ttfamily,numbersep=5pt]
cloudra/base/tools_class.py:97	
cloudra/base/tools_class.py:98	    def toJson(self) -> dict:
cloudra/base/tools_class.py:99	        '''实例序列化为 dict，用于日志或上下文回传。'''
cloudra/base/tools_class.py:100	        return self.model_dump(mode="json")
cloudra/base/tools_class.py:101	
cloudra/base/tools_class.py:102	
cloudra/base/tools_class.py:103	class output_schema(BaseModel):
cloudra/base/tools_class.py:104	    '''
cloudra/base/tools_class.py:105	    工具输出模式基类，定义工具输出的格式。
cloudra/base/tools_class.py:106	    子类按普通 Pydantic Model 声明字段即可。
cloudra/base/tools_class.py:107	    '''
cloudra/base/tools_class.py:108	
cloudra/base/tools_class.py:109	    def toJson(self) -> dict:
cloudra/base/tools_class.py:110	        return self.model_dump(mode="json")
cloudra/base/tools_class.py:111	
cloudra/base/tools_class.py:112	
cloudra/base/tools_class.py:113	class Tool(BaseModel):
cloudra/base/tools_class.py:114	    '''
cloudra/base/tools_class.py:115	    工具基类。所有具体工具应继承此类。
cloudra/base/tools_class.py:116	
cloudra/base/tools_class.py:117	    关于 input_schema / output_schema：
cloudra/base/tools_class.py:118	    - 这里存放的是具体的 schema 类（如 TextReader_input_schema），而非实例。
cloudra/base/tools_class.py:119	    - 使用 Any 类型并开启 arbitrary_types_allowed，避免 Pydantic 对类对象做字段校验。
cloudra/base/tools_class.py:120	    - 子模块中覆盖时必须带类型注解，例如：
cloudra/base/tools_class.py:121	          input_schema: Any = TextReader_input_schema
cloudra/base/tools_class.py:122	    '''
cloudra/base/tools_class.py:123	    model_config = {"arbitrary_types_allowed": True}
cloudra/base/tools_class.py:124	
cloudra/base/tools_class.py:125	    # 唯一标识符
cloudra/base/tools_class.py:126	    tool_id: str = ''
cloudra/base/tools_class.py:127	
cloudra/base/tools_class.py:128	    # 工具名称（即 OpenAI function.name）
cloudra/base/tools_class.py:129	    tool_name: str = ''
cloudra/base/tools_class.py:130	
cloudra/base/tools_class.py:131	    # 工具描述。改造后应仅保留【功能描述】+【使用场景】，不再包含参数细节。
cloudra/base/tools_class.py:132	    # 参数说明通过 input_schema 的 Field(description="...") 提供，自动生成到 parameters 中。
cloudra/base/tools_class.py:133	    tool_description: str = ''
cloudra/base/tools_class.py:134	
cloudra/base/tools_class.py:135	    # 工具的输入模式类（存放类对象，非实例）
cloudra/base/tools_class.py:136	    input_schema: Any = input_schema
cloudra/base/tools_class.py:137	
cloudra/base/tools_class.py:138	    # 工具的输出模式类（存放类对象，非实例）
cloudra/base/tools_class.py:139	    output_schema: Any = output_schema
cloudra/base/tools_class.py:140	
cloudra/base/tools_class.py:141	    # 引用调用该工具的 Agent 实例
cloudra/base/tools_class.py:142	    subagent: Any = None
cloudra/base/tools_class.py:143	
cloudra/base/tools_class.py:144	    def is_read(self) -> bool:
cloudra/base/tools_class.py:145	        return False
cloudra/base/tools_class.py:146	
cloudra/base/tools_class.py:147	    def is_write(self) -> bool:
cloudra/base/tools_class.py:148	        return False
\end{lstlisting}
\clearpage
\begin{lstlisting}[firstnumber=149,numbers=left,numberstyle=\tiny\ttfamily,numbersep=5pt]
cloudra/base/tools_class.py:149	
cloudra/base/tools_class.py:150	    def is_manage(self) -> bool:
cloudra/base/tools_class.py:151	        return False
cloudra/base/tools_class.py:152	
cloudra/base/tools_class.py:153	    def is_script(self) -> bool:
cloudra/base/tools_class.py:154	        return False
cloudra/base/tools_class.py:155	
cloudra/base/tools_class.py:156	    def is_destructive(self) -> bool:
cloudra/base/tools_class.py:157	        return False
cloudra/base/tools_class.py:158	
cloudra/base/tools_class.py:159	    def tool_permission_all(self) -> Permission:
cloudra/base/tools_class.py:160	        """Return Permission flags for this tool."""
cloudra/base/tools_class.py:161	        perms = Permission.NONE
cloudra/base/tools_class.py:162	        if self.is_read():
cloudra/base/tools_class.py:163	            perms |= Permission.READ
cloudra/base/tools_class.py:164	        if self.is_write():
cloudra/base/tools_class.py:165	            perms |= Permission.WRITE
cloudra/base/tools_class.py:166	        if self.is_manage():
cloudra/base/tools_class.py:167	            perms |= Permission.MANAGE
cloudra/base/tools_class.py:168	        if self.is_script():
cloudra/base/tools_class.py:169	            perms |= Permission.SCRIPT
cloudra/base/tools_class.py:170	        if self.is_destructive():
cloudra/base/tools_class.py:171	            perms |= Permission.DESTRUCTIVE
cloudra/base/tools_class.py:172	        return perms
cloudra/base/tools_class.py:173	
cloudra/base/tools_class.py:174	    async def tool_call(self, user_input: input_schema, session: Session) -> output_schema:
cloudra/base/tools_class.py:175	        '''
cloudra/base/tools_class.py:176	        调用工具。子类必须重写此方法。
cloudra/base/tools_class.py:177	        执行结果以 output_schema 实例返回。
cloudra/base/tools_class.py:178	        '''
cloudra/base/tools_class.py:179	        raise NotImplementedError("子类必须实现 tool_call 方法")
cloudra/base/tools_class.py:180	
cloudra/base/tools_class.py:181	    async def safe_call(self, user_input: input_schema, session: Session) -> output_schema:
cloudra/base/tools_class.py:182	        '''
cloudra/base/tools_class.py:183	        带有通用异常捕获的工具调用包装器。
cloudra/base/tools_class.py:184	        工具执行失败时抛出包含工具名的 RuntimeError，由 reAct.act_result() 统一捕获并回传给 LLM。
cloudra/base/tools_class.py:185	        消息以 [Tool error] 前缀 + 工具名 + 恢复指引成形；react 层识别该前缀后不再二次套壳。
cloudra/base/tools_class.py:186	        '''
cloudra/base/tools_class.py:187	        try:
cloudra/base/tools_class.py:188	            return await self.tool_call(user_input, session)
cloudra/base/tools_class.py:189	        except Exception as e:
cloudra/base/tools_class.py:190	            error_msg = (
cloudra/base/tools_class.py:191	                f"[Tool error] {self.tool_name}: {str(e)}. "
cloudra/base/tools_class.py:192	                "Fix the arguments or use a different tool, then retry."
cloudra/base/tools_class.py:193	            )
cloudra/base/tools_class.py:194	            raise RuntimeError(error_msg) from e
cloudra/base/tools_class.py:195	
cloudra/base/tools_class.py:196	    async def ensure_safety(
cloudra/base/tools_class.py:197	        self,
cloudra/base/tools_class.py:198	        session: Session,
cloudra/base/tools_class.py:199	        details: str,
cloudra/base/tools_class.py:200	        tool_args: dict | None = None,
\end{lstlisting}
\clearpage
\begin{lstlisting}[firstnumber=201,numbers=left,numberstyle=\tiny\ttfamily,numbersep=5pt]
cloudra/base/tools_class.py:201	    ) -> tuple[bool, str]:
cloudra/base/tools_class.py:202	        '''
cloudra/base/tools_class.py:203	        安全检查：如果工具被标记为 destructive，则必须经过审批。
cloudra/base/tools_class.py:204	        默认通过数据库等待用户审批；若启用了 YOLO 模式，则调用独立 Agent 自动判断。
cloudra/base/tools_class.py:205	        返回 (是否通过, 拒绝原因)。通过时拒绝原因为空字符串。
cloudra/base/tools_class.py:206	        '''
cloudra/base/tools_class.py:207	        from cloudra.infra.yolo_guard import YoloGuard, is_yolo_enabled
cloudra/base/tools_class.py:208	
cloudra/base/tools_class.py:209	        if is_yolo_enabled():
cloudra/base/tools_class.py:210	            guard = YoloGuard()
cloudra/base/tools_class.py:211	            return await self._yolo_judge(session, tool_args or {}, guard)
cloudra/base/tools_class.py:212	
cloudra/base/tools_class.py:213	        db_path = _tool_db_path(self.subagent)
cloudra/base/tools_class.py:214	        risk_summary = next(
cloudra/base/tools_class.py:215	            (
cloudra/base/tools_class.py:216	                line.partition(":")[2].strip()
cloudra/base/tools_class.py:217	                for line in details.splitlines()
cloudra/base/tools_class.py:218	                if line.lower().startswith("action:") or line.startswith("【动作】")
cloudra/base/tools_class.py:219	            ),
cloudra/base/tools_class.py:220	            details.splitlines()[0] if details.splitlines() else "Dangerous tool call",
cloudra/base/tools_class.py:221	        )
cloudra/base/tools_class.py:222	        candidate_diff = ""
cloudra/base/tools_class.py:223	        for marker in ("Unified Diff:\n", "【Unified Diff】:\n"):
cloudra/base/tools_class.py:224	            if marker in details:
cloudra/base/tools_class.py:225	                candidate_diff = details.partition(marker)[2]
cloudra/base/tools_class.py:226	                break
cloudra/base/tools_class.py:227	
cloudra/base/tools_class.py:228	        return await self.request_approval(
cloudra/base/tools_class.py:229	            db_path,
cloudra/base/tools_class.py:230	            session.id,
cloudra/base/tools_class.py:231	            details,
cloudra/base/tools_class.py:232	            tool_name=self.tool_name,
cloudra/base/tools_class.py:233	            normalized_arguments=tool_args or {},
cloudra/base/tools_class.py:234	            risk_summary=risk_summary,
cloudra/base/tools_class.py:235	            candidate_diff=candidate_diff,
cloudra/base/tools_class.py:236	        )
cloudra/base/tools_class.py:237	
cloudra/base/tools_class.py:238	    async def record_session_event(
cloudra/base/tools_class.py:239	        self,
cloudra/base/tools_class.py:240	        session: Session,
cloudra/base/tools_class.py:241	        event_type: str,
cloudra/base/tools_class.py:242	        data: dict,
cloudra/base/tools_class.py:243	    ) -> None:
cloudra/base/tools_class.py:244	        """尽力记录工具事件；遥测失败不能反向破坏工具主流程。"""
cloudra/base/tools_class.py:245	        from cloudra.services import session_store as _ss
cloudra/base/tools_class.py:246	
cloudra/base/tools_class.py:247	        try:
cloudra/base/tools_class.py:248	            db = ConversationDB.from_path(_tool_db_path(self.subagent))
cloudra/base/tools_class.py:249	            await _ss.arecord_event(db, session.id, event_type, data)
cloudra/base/tools_class.py:250	        except Exception:
cloudra/base/tools_class.py:251	            pass
cloudra/base/tools_class.py:252	
\end{lstlisting}
\clearpage
\begin{lstlisting}[firstnumber=253,numbers=left,numberstyle=\tiny\ttfamily,numbersep=5pt]
cloudra/base/tools_class.py:253	    async def _yolo_judge(
cloudra/base/tools_class.py:254	        self,
cloudra/base/tools_class.py:255	        session: Session,
cloudra/base/tools_class.py:256	        tool_args: dict,
cloudra/base/tools_class.py:257	        guard: 'YoloGuard',
cloudra/base/tools_class.py:258	    ) -> tuple[bool, str]:
cloudra/base/tools_class.py:259	        """调用 YoloGuard 做自动判断。"""
cloudra/base/tools_class.py:260	        approved, reason = await guard.judge(
cloudra/base/tools_class.py:261	            self.tool_name, tool_args, self.subagent, session
cloudra/base/tools_class.py:262	        )
cloudra/base/tools_class.py:263	        if approved:
cloudra/base/tools_class.py:264	            print("\033[92m[安全审计/YOLO]\033[0m 自动审批通过。")
cloudra/base/tools_class.py:265	            return True, ""
cloudra/base/tools_class.py:266	        else:
cloudra/base/tools_class.py:267	            reason_text = reason or "YOLO Agent 认为该操作存在风险"
cloudra/base/tools_class.py:268	            print(f"\033[91m[安全审计/YOLO]\033[0m 自动审批拒绝。原因: {reason_text}")
cloudra/base/tools_class.py:269	            return False, reason_text
cloudra/base/tools_class.py:270	
cloudra/base/tools_class.py:271	    async def request_approval(
cloudra/base/tools_class.py:272	        self,
cloudra/base/tools_class.py:273	        db_path: str,
cloudra/base/tools_class.py:274	        session_id: str,
cloudra/base/tools_class.py:275	        details: str,
cloudra/base/tools_class.py:276	        *,
cloudra/base/tools_class.py:277	        tool_name: str = "",
cloudra/base/tools_class.py:278	        normalized_arguments: dict | None = None,
cloudra/base/tools_class.py:279	        risk_summary: str = "",
cloudra/base/tools_class.py:280	        candidate_diff: str = "",
cloudra/base/tools_class.py:281	    ) -> tuple[bool, str]:
cloudra/base/tools_class.py:282	        handler = _APPROVAL_HANDLER.get()
cloudra/base/tools_class.py:283	        if handler is None:
cloudra/base/tools_class.py:284	            # v1 时代"SQLite + 事件总线等待"的审批回退已随旧服务层删除。
cloudra/base/tools_class.py:285	            raise RuntimeError(
cloudra/base/tools_class.py:286	                "no approval handler installed: runtime_v2 LegacyToolExecutor always sets one"
cloudra/base/tools_class.py:287	            )
cloudra/base/tools_class.py:288	        # runtime_v2 上下文：审批请求与等待由内核 broker 承担，
cloudra/base/tools_class.py:289	        # 不碰 conversation SQLite。
cloudra/base/tools_class.py:290	        return await handler(
cloudra/base/tools_class.py:291	            details=details,
cloudra/base/tools_class.py:292	            tool_name=tool_name,
cloudra/base/tools_class.py:293	            normalized_arguments=normalized_arguments or {},
cloudra/base/tools_class.py:294	            risk_summary=risk_summary,
cloudra/base/tools_class.py:295	            candidate_diff=candidate_diff,
cloudra/base/tools_class.py:296	        )
cloudra/base/tools_class.py:297	
cloudra/base/tools_class.py:298	    @classmethod
cloudra/base/tools_class.py:299	    def _get_field_default(cls, field_name: str):
cloudra/base/tools_class.py:300	        '''获取类字段的默认值，兼容 Pydantic v2（类级别无法直接访问普通字段）。'''
cloudra/base/tools_class.py:301	        field = cls.model_fields.get(field_name)
cloudra/base/tools_class.py:302	        if field and field.default is not PydanticUndefined:
cloudra/base/tools_class.py:303	            return field.default
cloudra/base/tools_class.py:304	        return None
\end{lstlisting}
\clearpage
\begin{lstlisting}[firstnumber=305,numbers=left,numberstyle=\tiny\ttfamily,numbersep=5pt]
cloudra/base/tools_class.py:305	
cloudra/base/tools_class.py:306	    def tool_schema_json(self) -> dict:
cloudra/base/tools_class.py:307	        '''
cloudra/base/tools_class.py:308	        返回标准 OpenAI function 格式，可直接用于 client.chat.completions.create(..., tools=[...])。
cloudra/base/tools_class.py:309	        优先读取实例级别的 tool_description，使动态生成的描述能够生效。
cloudra/base/tools_class.py:310	        '''
cloudra/base/tools_class.py:311	        # HACK: 这里用 getattr + or fallback 链是因为 Pydantic v2 的类字段默认值在实例上可能
cloudra/base/tools_class.py:312	        # 被 shadow，但某些情况下实例属性又还没覆盖。动态 description（如 Bashtools 的权限列表、
cloudra/base/tools_class.py:313	        # GetSkill 的 skills 列表）必须在 __init__ 之后才能确定，因此必须是实例方法而非类方法。
cloudra/base/tools_class.py:314	        # 如果后续有更好方式（如 descriptor / property），可替换掉这个 ugly fallback。
cloudra/base/tools_class.py:315	        description = getattr(self, 'tool_description', '') or self._get_field_default('tool_description') or ''
cloudra/base/tools_class.py:316	        return {
cloudra/base/tools_class.py:317	            "type": "function",
cloudra/base/tools_class.py:318	            "function": {
cloudra/base/tools_class.py:319	                "name": getattr(self, 'tool_name', '') or self._get_field_default('tool_name') or '',
cloudra/base/tools_class.py:320	                "description": description,
cloudra/base/tools_class.py:321	                "parameters": (self._get_field_default('input_schema') or input_schema).to_parameters_schema(),
cloudra/base/tools_class.py:322	            }
cloudra/base/tools_class.py:323	        }
cloudra/base/tools_class.py:324	
cloudra/base/tools_class.py:325	    def toJson(self) -> dict:
cloudra/base/tools_class.py:326	        '''
cloudra/base/tools_class.py:327	        实例级序列化，主要用于日志或调试。
cloudra/base/tools_class.py:328	        若需获取 OpenAI schema，请使用类方法 tool_schema_json()。
cloudra/base/tools_class.py:329	        '''
cloudra/base/tools_class.py:330	        return {
cloudra/base/tools_class.py:331	            "tool_id": self.tool_id,
cloudra/base/tools_class.py:332	            "tool_name": self.tool_name,
cloudra/base/tools_class.py:333	            "tool_description": self.tool_description,
cloudra/base/tools_class.py:334	        }
cloudra/base/tools_class.py:335	
cloudra/base/tools_class.py:336	    @classmethod
cloudra/base/tools_class.py:337	    def build_input(cls, arg_dict: dict) -> input_schema:
cloudra/base/tools_class.py:338	        '''
cloudra/base/tools_class.py:339	        将 LLM 返回的参数字典转换为 input_schema 实例。
cloudra/base/tools_class.py:340	        供 reAct.py 调用，替代原有的 tool.input_schema.toClass(tool_args)。
cloudra/base/tools_class.py:341	        '''
cloudra/base/tools_class.py:342	        schema_cls = cls._get_field_default('input_schema') or input_schema
cloudra/base/tools_class.py:343	        return schema_cls.from_args(arg_dict)
cloudra/benchmark/__init__.py:1	"""Benchmark module: benchmark task schema, task generators, and the scoring pipeline.
cloudra/benchmark/__init__.py:2	
cloudra/benchmark/__init__.py:3	Public surface:
cloudra/benchmark/__init__.py:4	- `schema`: TaskSpec / GoldSpec / RubricSpec / mechanism-graph models.
cloudra/benchmark/__init__.py:5	- `gen_easy_tasks`: generator turning ideapool/reaction_candidates/easy.md
cloudra/benchmark/__init__.py:6	  into ideapool/tasks/easy/*.json.
cloudra/benchmark/__init__.py:7	- `graph_similarity`: chemical-parameterized GED between two mechanism graphs.
cloudra/benchmark/__init__.py:8	- `integrity`: sandbox forensic scan + evidence binding (Layer 1).
cloudra/benchmark/__init__.py:9	- `selfconsistency`: physical self-consistency checks (Layer 2).
cloudra/benchmark/__init__.py:10	- `accuracy`: Gaussian-tolerance barrier scoring (Layer 3).
cloudra/benchmark/__init__.py:11	- `method_policy`: method-whitelist interception.
cloudra/benchmark/__init__.py:12	- `runner`: independent benchmark runner with its own data directory.
cloudra/benchmark/__init__.py:13	"""
\end{lstlisting}
\clearpage
\begin{lstlisting}[firstnumber=14,numbers=left,numberstyle=\tiny\ttfamily,numbersep=5pt]
cloudra/benchmark/__init__.py:14	
cloudra/benchmark/__init__.py:15	from cloudra.benchmark.schema import (  # noqa: F401
cloudra/benchmark/__init__.py:16	    AccuracyRubric,
cloudra/benchmark/__init__.py:17	    BondChangeSet,
cloudra/benchmark/__init__.py:18	    ExpectedOutcome,
cloudra/benchmark/__init__.py:19	    GoldMechanismGraph,
cloudra/benchmark/__init__.py:20	    GoldSpec,
cloudra/benchmark/__init__.py:21	    GraphSimilarityRubric,
cloudra/benchmark/__init__.py:22	    ImpossibleCaseRubric,
cloudra/benchmark/__init__.py:23	    IntegrityRubric,
cloudra/benchmark/__init__.py:24	    MechanismEdge,
cloudra/benchmark/__init__.py:25	    MechanismNode,
cloudra/benchmark/__init__.py:26	    MethodPolicy,
cloudra/benchmark/__init__.py:27	    NodeRole,
cloudra/benchmark/__init__.py:28	    PromptSpec,
cloudra/benchmark/__init__.py:29	    ReferenceComputation,
cloudra/benchmark/__init__.py:30	    RubricSpec,
cloudra/benchmark/__init__.py:31	    SelfConsistencyRubric,
cloudra/benchmark/__init__.py:32	    SystemSpec,
cloudra/benchmark/__init__.py:33	    TaskSpec,
cloudra/benchmark/__init__.py:34	    Tier,
cloudra/benchmark/__init__.py:35	    VerificationStatus,
cloudra/benchmark/__init__.py:36	)
cloudra/benchmark/__main__.py:1	"""python -m cloudra.benchmark — run the benchmark task bank.
cloudra/benchmark/__main__.py:2	
cloudra/benchmark/__main__.py:3	Default is a dry-run (no-op backend: sandboxes stay empty, every task scores
cloudra/benchmark/__main__.py:4	0 through the real scoring pipeline — useful for checking the harness itself).
cloudra/benchmark/__main__.py:5	`--backend runtime_v2` runs each task on the real in-process runtime_v2 kernel
cloudra/benchmark/__main__.py:6	with the method whitelist enforced.
cloudra/benchmark/__main__.py:7	
cloudra/benchmark/__main__.py:8	Example:
cloudra/benchmark/__main__.py:9	    python -m cloudra.benchmark --backend runtime_v2 --data-dir bench_out --only C1,D3
cloudra/benchmark/__main__.py:10	"""
cloudra/benchmark/__main__.py:11	
cloudra/benchmark/__main__.py:12	from __future__ import annotations
cloudra/benchmark/__main__.py:13	
cloudra/benchmark/__main__.py:14	import argparse
cloudra/benchmark/__main__.py:15	import os
cloudra/benchmark/__main__.py:16	import sys
cloudra/benchmark/__main__.py:17	from pathlib import Path
cloudra/benchmark/__main__.py:18	
cloudra/benchmark/__main__.py:19	
cloudra/benchmark/__main__.py:20	def main(argv: list[str] | None = None) -> int:
cloudra/benchmark/__main__.py:21	    from cloudra.benchmark.runner import BenchmarkRunner, load_tasks
cloudra/benchmark/__main__.py:22	
cloudra/benchmark/__main__.py:23	    ap = argparse.ArgumentParser(prog="cloudra.benchmark")
cloudra/benchmark/__main__.py:24	    ap.add_argument("--tasks-dir", type=Path, default=Path("ideapool/tasks/easy"))
cloudra/benchmark/__main__.py:25	    ap.add_argument("--data-dir", type=Path, default=Path("benchmark_data"),
cloudra/benchmark/__main__.py:26	                    help="benchmark data dir (sandboxes + benchmark.db)")
cloudra/benchmark/__main__.py:27	    ap.add_argument("--backend", choices=["none", "runtime_v2"], default="none")
cloudra/benchmark/__main__.py:28	    ap.add_argument("--only", default="", help="comma-separated task ids (default: all)")
cloudra/benchmark/__main__.py:29	    ap.add_argument("--timeout", type=float, default=7200.0, help="per-task run timeout, seconds")
\end{lstlisting}
\clearpage
\begin{lstlisting}[firstnumber=30,numbers=left,numberstyle=\tiny\ttfamily,numbersep=5pt]
cloudra/benchmark/__main__.py:30	    ap.add_argument("--concurrency", type=int, default=1,
cloudra/benchmark/__main__.py:31	                    help="parallel tasks (threads, per-task event loop + sandbox; default 1 = serial)")
cloudra/benchmark/__main__.py:32	    ap.add_argument("--api-pool", type=Path, default=None,
cloudra/benchmark/__main__.py:33	                    help='JSON endpoint pool for anti-429: [{"base_url","api_key","label"}, ...]; '
cloudra/benchmark/__main__.py:34	                         "each LLM call picks a random endpoint, 429 fails over to the next")
cloudra/benchmark/__main__.py:35	    args = ap.parse_args(argv)
cloudra/benchmark/__main__.py:36	
cloudra/benchmark/__main__.py:37	    tasks = load_tasks(args.tasks_dir)
cloudra/benchmark/__main__.py:38	    if args.only:
cloudra/benchmark/__main__.py:39	        wanted = {t.strip() for t in args.only.split(",") if t.strip()}
cloudra/benchmark/__main__.py:40	        tasks = [t for t in tasks if t.task_id in wanted]
cloudra/benchmark/__main__.py:41	    if not tasks:
cloudra/benchmark/__main__.py:42	        print("no tasks matched", file=sys.stderr)
cloudra/benchmark/__main__.py:43	        return 2
cloudra/benchmark/__main__.py:44	
cloudra/benchmark/__main__.py:45	    backend = None
cloudra/benchmark/__main__.py:46	    if args.backend == "runtime_v2":
cloudra/benchmark/__main__.py:47	        # 运行时数据隔离必须在首次读取全局配置前完成；config_file 独立于
cloudra/benchmark/__main__.py:48	        # data_dir，LLM 凭据不受影响。YAML 里的 data_dir 键会覆盖环境变量，
cloudra/benchmark/__main__.py:49	        # 所以两处都钉死。
cloudra/benchmark/__main__.py:50	        isolated = (args.data_dir / "runtime").resolve()
cloudra/benchmark/__main__.py:51	        os.environ["CLOUDRA_DATA_DIR"] = str(isolated)
cloudra/benchmark/__main__.py:52	        from cloudra.infra.config import get_config
cloudra/benchmark/__main__.py:53	
cloudra/benchmark/__main__.py:54	        get_config().data_dir = isolated
cloudra/benchmark/__main__.py:55	        llm = None
cloudra/benchmark/__main__.py:56	        if args.api_pool is not None:
cloudra/benchmark/__main__.py:57	            from cloudra.benchmark.api_pool import PooledLLMProvider, load_endpoints
cloudra/benchmark/__main__.py:58	
cloudra/benchmark/__main__.py:59	            llm = PooledLLMProvider(load_endpoints(args.api_pool))
cloudra/benchmark/__main__.py:60	        from cloudra.benchmark.backend import RuntimeV2Backend
cloudra/benchmark/__main__.py:61	
cloudra/benchmark/__main__.py:62	        backend = RuntimeV2Backend(run_timeout_s=args.timeout, llm=llm)
cloudra/benchmark/__main__.py:63	
cloudra/benchmark/__main__.py:64	    runner = BenchmarkRunner(args.data_dir, backend=backend)
cloudra/benchmark/__main__.py:65	    results = runner.run_all(tasks, concurrency=max(1, args.concurrency))
cloudra/benchmark/__main__.py:66	    for r in results:
cloudra/benchmark/__main__.py:67	        print(
cloudra/benchmark/__main__.py:68	            f"{r.task_id:>6}  final={r.final_score:.3f}  "
cloudra/benchmark/__main__.py:69	            f"integrity={int(r.integrity_passed)} graph={r.graph_score:.3f} "
cloudra/benchmark/__main__.py:70	            f"selfcons={int(r.selfconsistency_passed)} accuracy={r.accuracy_score:.3f}"
cloudra/benchmark/__main__.py:71	        )
cloudra/benchmark/__main__.py:72	    return 0
cloudra/benchmark/__main__.py:73	
cloudra/benchmark/__main__.py:74	
cloudra/benchmark/__main__.py:75	if __name__ == "__main__":
cloudra/benchmark/__main__.py:76	    sys.exit(main())
cloudra/benchmark/accuracy.py:1	"""Layer 3 — accuracy: Gaussian-tolerance barrier scoring against the gold
cloudra/benchmark/accuracy.py:2	reference. Per-task sigma; a configurable cap on how much accuracy can cost.
cloudra/benchmark/accuracy.py:3	"""
cloudra/benchmark/accuracy.py:4	
cloudra/benchmark/accuracy.py:5	from __future__ import annotations
\end{lstlisting}
\clearpage
\begin{lstlisting}[firstnumber=6,numbers=left,numberstyle=\tiny\ttfamily,numbersep=5pt]
cloudra/benchmark/accuracy.py:6	
cloudra/benchmark/accuracy.py:7	import math
cloudra/benchmark/accuracy.py:8	from dataclasses import dataclass, field
cloudra/benchmark/accuracy.py:9	
cloudra/benchmark/accuracy.py:10	from cloudra.benchmark.schema import AccuracyRubric, ReferenceComputation
cloudra/benchmark/accuracy.py:11	
cloudra/benchmark/accuracy.py:12	
cloudra/benchmark/accuracy.py:13	@dataclass
cloudra/benchmark/accuracy.py:14	class AccuracyReport:
cloudra/benchmark/accuracy.py:15	    enabled: bool = True
cloudra/benchmark/accuracy.py:16	    agent_barrier_kcal: float | None = None
cloudra/benchmark/accuracy.py:17	    gold_barrier_kcal: float | None = None
cloudra/benchmark/accuracy.py:18	    sigma_kcal: float = 3.0
cloudra/benchmark/accuracy.py:19	    gaussian_score: float = 0.0          # in [0, 1], 1 = exact match
cloudra/benchmark/accuracy.py:20	    capped_penalty: float = 0.0          # penalty actually applied after cap
cloudra/benchmark/accuracy.py:21	    final_score: float = 0.0             # 1 - capped_penalty
cloudra/benchmark/accuracy.py:22	    notes: list[str] = field(default_factory=list)
cloudra/benchmark/accuracy.py:23	
cloudra/benchmark/accuracy.py:24	
cloudra/benchmark/accuracy.py:25	def gaussian_barrier_score(agent: float, gold: float, sigma: float) -> float:
cloudra/benchmark/accuracy.py:26	    """exp(-(agent-gold)^2 / (2 sigma^2)) in [0, 1]."""
cloudra/benchmark/accuracy.py:27	    if sigma <= 0:
cloudra/benchmark/accuracy.py:28	        return 1.0 if agent == gold else 0.0
cloudra/benchmark/accuracy.py:29	    d = agent - gold
cloudra/benchmark/accuracy.py:30	    return math.exp(-(d * d) / (2.0 * sigma * sigma))
cloudra/benchmark/accuracy.py:31	
cloudra/benchmark/accuracy.py:32	
cloudra/benchmark/accuracy.py:33	def score_accuracy(
cloudra/benchmark/accuracy.py:34	    agent_barrier_kcal: float | None,
cloudra/benchmark/accuracy.py:35	    gold: ReferenceComputation,
cloudra/benchmark/accuracy.py:36	    rubric: AccuracyRubric | None = None,
cloudra/benchmark/accuracy.py:37	) -> AccuracyReport:
cloudra/benchmark/accuracy.py:38	    rubric = rubric or AccuracyRubric()
cloudra/benchmark/accuracy.py:39	    rep = AccuracyReport(enabled=rubric.enabled, sigma_kcal=rubric.sigma_kcal_mol)
cloudra/benchmark/accuracy.py:40	    rep.agent_barrier_kcal = agent_barrier_kcal
cloudra/benchmark/accuracy.py:41	    rep.gold_barrier_kcal = gold.barrier_kcal_mol
cloudra/benchmark/accuracy.py:42	
cloudra/benchmark/accuracy.py:43	    if not rubric.enabled:
cloudra/benchmark/accuracy.py:44	        rep.gaussian_score = 1.0
cloudra/benchmark/accuracy.py:45	        rep.final_score = 1.0
cloudra/benchmark/accuracy.py:46	        rep.notes.append("accuracy disabled for this task; scored on self-consistency only")
cloudra/benchmark/accuracy.py:47	        return rep
cloudra/benchmark/accuracy.py:48	    if gold.barrier_kcal_mol is None:
cloudra/benchmark/accuracy.py:49	        rep.gaussian_score = 1.0
cloudra/benchmark/accuracy.py:50	        rep.final_score = 1.0
cloudra/benchmark/accuracy.py:51	        rep.notes.append("no gold barrier; accuracy not scored")
cloudra/benchmark/accuracy.py:52	        return rep
cloudra/benchmark/accuracy.py:53	    if agent_barrier_kcal is None:
cloudra/benchmark/accuracy.py:54	        rep.gaussian_score = 0.0
cloudra/benchmark/accuracy.py:55	        rep.capped_penalty = rubric.max_penalty_fraction
cloudra/benchmark/accuracy.py:56	        rep.final_score = 1.0 - rep.capped_penalty
cloudra/benchmark/accuracy.py:57	        rep.notes.append("agent barrier missing")
\end{lstlisting}
\clearpage
\begin{lstlisting}[firstnumber=58,numbers=left,numberstyle=\tiny\ttfamily,numbersep=5pt]
cloudra/benchmark/accuracy.py:58	        return rep
cloudra/benchmark/accuracy.py:59	
cloudra/benchmark/accuracy.py:60	    g = gaussian_barrier_score(agent_barrier_kcal, gold.barrier_kcal_mol, rubric.sigma_kcal_mol)
cloudra/benchmark/accuracy.py:61	    rep.gaussian_score = g
cloudra/benchmark/accuracy.py:62	    raw_penalty = 1.0 - g
cloudra/benchmark/accuracy.py:63	    rep.capped_penalty = min(raw_penalty, rubric.max_penalty_fraction)
cloudra/benchmark/accuracy.py:64	    rep.final_score = 1.0 - rep.capped_penalty
cloudra/benchmark/accuracy.py:65	    if raw_penalty > rubric.max_penalty_fraction:
cloudra/benchmark/accuracy.py:66	        rep.notes.append(
cloudra/benchmark/accuracy.py:67	            f"raw penalty {raw_penalty:.3f} capped at {rubric.max_penalty_fraction}"
cloudra/benchmark/accuracy.py:68	        )
cloudra/benchmark/accuracy.py:69	    return rep
cloudra/benchmark/api_pool.py:1	"""API endpoint pool for benchmark runs.
cloudra/benchmark/api_pool.py:2	
cloudra/benchmark/api_pool.py:3	Anti-429 for parallel benchmark execution: every LLM call picks a random
cloudra/benchmark/api_pool.py:4	endpoint from the pool (spreading load), and if the call still exhausts the
cloudra/benchmark/api_pool.py:5	client-level backoff retries with a rate-limit error, the next endpoint takes
cloudra/benchmark/api_pool.py:6	over (no repeats within one call).
cloudra/benchmark/api_pool.py:7	
cloudra/benchmark/api_pool.py:8	Scope by design: pool members share the *model* from the agent's settings —
cloudra/benchmark/api_pool.py:9	only the transport (base_url / api_key) is pinned per endpoint. Empty fields
cloudra/benchmark/api_pool.py:10	fall back to the agent's settings, so a pool can be just a list of extra keys
cloudra/benchmark/api_pool.py:11	for the same gateway. Non-429 errors are not retried here: they fail the turn
cloudra/benchmark/api_pool.py:12	(the runner's own turn-retry handles transient crashes as before).
cloudra/benchmark/api_pool.py:13	"""
cloudra/benchmark/api_pool.py:14	
cloudra/benchmark/api_pool.py:15	from __future__ import annotations
cloudra/benchmark/api_pool.py:16	
cloudra/benchmark/api_pool.py:17	import contextvars
cloudra/benchmark/api_pool.py:18	import json
cloudra/benchmark/api_pool.py:19	import logging
cloudra/benchmark/api_pool.py:20	import random
cloudra/benchmark/api_pool.py:21	from dataclasses import dataclass
cloudra/benchmark/api_pool.py:22	from pathlib import Path
cloudra/benchmark/api_pool.py:23	from typing import Any, Optional
cloudra/benchmark/api_pool.py:24	
cloudra/benchmark/api_pool.py:25	from cloudra.runtime_v2.adapters import LegacyLLMProvider
cloudra/benchmark/api_pool.py:26	
cloudra/benchmark/api_pool.py:27	logger = logging.getLogger(__name__)
cloudra/benchmark/api_pool.py:28	
cloudra/benchmark/api_pool.py:29	
cloudra/benchmark/api_pool.py:30	@dataclass(frozen=True)
cloudra/benchmark/api_pool.py:31	class Endpoint:
cloudra/benchmark/api_pool.py:32	    """One pool member; empty fields fall back to the agent's settings."""
cloudra/benchmark/api_pool.py:33	
cloudra/benchmark/api_pool.py:34	    base_url: str = ""
cloudra/benchmark/api_pool.py:35	    api_key: str = ""
cloudra/benchmark/api_pool.py:36	    label: str = ""  # 日志用；空则取 base_url
cloudra/benchmark/api_pool.py:37	
cloudra/benchmark/api_pool.py:38	    def display(self) -> str:
cloudra/benchmark/api_pool.py:39	        return self.label or self.base_url or "(settings default)"
cloudra/benchmark/api_pool.py:40	
\end{lstlisting}
\clearpage
\begin{lstlisting}[firstnumber=41,numbers=left,numberstyle=\tiny\ttfamily,numbersep=5pt]
cloudra/benchmark/api_pool.py:41	
cloudra/benchmark/api_pool.py:42	# 当前调用钉住的端点。contextvars 按 asyncio task / 线程隔离——并行跑题时
cloudra/benchmark/api_pool.py:43	# 共享同一个 provider 实例也不会串端点。
cloudra/benchmark/api_pool.py:44	_current_endpoint: contextvars.ContextVar[Optional[Endpoint]] = contextvars.ContextVar(
cloudra/benchmark/api_pool.py:45	    "benchmark_api_pool_endpoint", default=None
cloudra/benchmark/api_pool.py:46	)
cloudra/benchmark/api_pool.py:47	
cloudra/benchmark/api_pool.py:48	
cloudra/benchmark/api_pool.py:49	def load_endpoints(path: Path | str) -> list[Endpoint]:
cloudra/benchmark/api_pool.py:50	    """Load a pool JSON: [{"base_url": ..., "api_key": ..., "label": ...}, ...]."""
cloudra/benchmark/api_pool.py:51	    raw = json.loads(Path(path).read_text(encoding="utf-8"))
cloudra/benchmark/api_pool.py:52	    if not isinstance(raw, list) or not raw:
cloudra/benchmark/api_pool.py:53	        raise ValueError(f"api pool file must be a non-empty JSON list: {path}")
cloudra/benchmark/api_pool.py:54	    endpoints = []
cloudra/benchmark/api_pool.py:55	    for i, item in enumerate(raw):
cloudra/benchmark/api_pool.py:56	        if not isinstance(item, dict):
cloudra/benchmark/api_pool.py:57	            raise ValueError(f"api pool entry #{i} must be an object: {item!r}")
cloudra/benchmark/api_pool.py:58	        endpoints.append(
cloudra/benchmark/api_pool.py:59	            Endpoint(
cloudra/benchmark/api_pool.py:60	                base_url=str(item.get("base_url") or ""),
cloudra/benchmark/api_pool.py:61	                api_key=str(item.get("api_key") or ""),
cloudra/benchmark/api_pool.py:62	                label=str(item.get("label") or ""),
cloudra/benchmark/api_pool.py:63	            )
cloudra/benchmark/api_pool.py:64	        )
cloudra/benchmark/api_pool.py:65	    return endpoints
cloudra/benchmark/api_pool.py:66	
cloudra/benchmark/api_pool.py:67	
cloudra/benchmark/api_pool.py:68	def _is_rate_limit(exc: BaseException) -> bool:
cloudra/benchmark/api_pool.py:69	    try:
cloudra/benchmark/api_pool.py:70	        import openai
cloudra/benchmark/api_pool.py:71	    except Exception:  # pragma: no cover - openai 是核心依赖
cloudra/benchmark/api_pool.py:72	        return False
cloudra/benchmark/api_pool.py:73	    return isinstance(exc, openai.RateLimitError)
cloudra/benchmark/api_pool.py:74	
cloudra/benchmark/api_pool.py:75	
cloudra/benchmark/api_pool.py:76	class PooledLLMProvider(LegacyLLMProvider):
cloudra/benchmark/api_pool.py:77	    """LegacyLLMProvider whose transport endpoint is pinned per call from a pool."""
cloudra/benchmark/api_pool.py:78	
cloudra/benchmark/api_pool.py:79	    def __init__(self, endpoints: list[Endpoint], *, rng: Any = None) -> None:
cloudra/benchmark/api_pool.py:80	        super().__init__()
cloudra/benchmark/api_pool.py:81	        self._endpoints = list(endpoints)
cloudra/benchmark/api_pool.py:82	        # rng 可注入（测试用固定顺序）；生产默认 Random()
cloudra/benchmark/api_pool.py:83	        self._rng = rng if rng is not None else random.Random()
cloudra/benchmark/api_pool.py:84	
cloudra/benchmark/api_pool.py:85	    def _make_client(self, settings):
cloudra/benchmark/api_pool.py:86	        ep = _current_endpoint.get()
cloudra/benchmark/api_pool.py:87	        if ep is None:
cloudra/benchmark/api_pool.py:88	            return super()._make_client(settings)
cloudra/benchmark/api_pool.py:89	        from cloudra.infra.llm_client import LLMClient
cloudra/benchmark/api_pool.py:90	
cloudra/benchmark/api_pool.py:91	        return LLMClient(
cloudra/benchmark/api_pool.py:92	            api_key=ep.api_key or settings.api_key,
\end{lstlisting}
\clearpage
\begin{lstlisting}[firstnumber=93,numbers=left,numberstyle=\tiny\ttfamily,numbersep=5pt]
cloudra/benchmark/api_pool.py:93	            base_url=ep.base_url or settings.base_url,
cloudra/benchmark/api_pool.py:94	            timeout=settings.timeout,
cloudra/benchmark/api_pool.py:95	        )
cloudra/benchmark/api_pool.py:96	
cloudra/benchmark/api_pool.py:97	    async def complete(self, agent, messages, tools, on_delta=None):
cloudra/benchmark/api_pool.py:98	        if not self._endpoints:
cloudra/benchmark/api_pool.py:99	            return await super().complete(agent, messages, tools, on_delta)
cloudra/benchmark/api_pool.py:100	        # 每次调用随机洗牌端点顺序：既摊初始负载，又保证 429 时不重复试同一个
cloudra/benchmark/api_pool.py:101	        order = self._rng.sample(self._endpoints, k=len(self._endpoints))
cloudra/benchmark/api_pool.py:102	        last_exc: Optional[BaseException] = None
cloudra/benchmark/api_pool.py:103	        for ep in order:
cloudra/benchmark/api_pool.py:104	            token = _current_endpoint.set(ep)
cloudra/benchmark/api_pool.py:105	            try:
cloudra/benchmark/api_pool.py:106	                return await super().complete(agent, messages, tools, on_delta)
cloudra/benchmark/api_pool.py:107	            except BaseException as exc:
cloudra/benchmark/api_pool.py:108	                if not _is_rate_limit(exc):
cloudra/benchmark/api_pool.py:109	                    raise
cloudra/benchmark/api_pool.py:110	                last_exc = exc
cloudra/benchmark/api_pool.py:111	                logger.warning("endpoint %s rate-limited (429); trying next", ep.display())
cloudra/benchmark/api_pool.py:112	            finally:
cloudra/benchmark/api_pool.py:113	                _current_endpoint.reset(token)
cloudra/benchmark/api_pool.py:114	        assert last_exc is not None
cloudra/benchmark/api_pool.py:115	        raise last_exc
cloudra/benchmark/backend.py:1	"""runtime_v2-backed agent backend for the benchmark runner.
cloudra/benchmark/backend.py:2	
cloudra/benchmark/backend.py:3	Per task this spins up an in-process runtime_v2 kernel (RuntimeCore +
cloudra/benchmark/backend.py:4	TurnPipeline + AgentRunner), spawns the task's coordinator agent with the task
cloudra/benchmark/backend.py:5	query against the per-task sandbox as the project workspace, waits for a
cloudra/benchmark/backend.py:6	terminal state, then extracts the agent-reported barrier from its final
cloudra/benchmark/backend.py:7	message. The method whitelist is enforced by appending `MethodPolicyStage` to
cloudra/benchmark/backend.py:8	the pipeline; compute jobs submitted via `submit_compute_job` are claimed and
cloudra/benchmark/backend.py:9	executed by a per-task in-process `SchedulerDaemon` (production runs it as a
cloudra/benchmark/backend.py:10	separate `cloudra scheduler` process instead).
cloudra/benchmark/backend.py:11	
cloudra/benchmark/backend.py:12	Process-wide isolation is the caller's job: the benchmark CLI points the
cloudra/benchmark/backend.py:13	global config's `data_dir` at the benchmark data dir before any tool runs, so
cloudra/benchmark/backend.py:14	session stores, compute jobs and MEMORY.md never touch the live conversation
cloudra/benchmark/backend.py:15	data dir. LLM credentials still come from the normal global config slots
cloudra/benchmark/backend.py:16	(coordinator/follower), unchanged.
cloudra/benchmark/backend.py:17	"""
cloudra/benchmark/backend.py:18	
cloudra/benchmark/backend.py:19	from __future__ import annotations
cloudra/benchmark/backend.py:20	
cloudra/benchmark/backend.py:21	import asyncio
cloudra/benchmark/backend.py:22	import logging
cloudra/benchmark/backend.py:23	import os
cloudra/benchmark/backend.py:24	import re
cloudra/benchmark/backend.py:25	import threading
cloudra/benchmark/backend.py:26	from pathlib import Path
cloudra/benchmark/backend.py:27	from typing import Any, Awaitable, Callable, Optional
cloudra/benchmark/backend.py:28	
cloudra/benchmark/backend.py:29	from cloudra.benchmark.policy_stage import MethodPolicyStage
\end{lstlisting}
\clearpage
\begin{lstlisting}[firstnumber=30,numbers=left,numberstyle=\tiny\ttfamily,numbersep=5pt]
cloudra/benchmark/backend.py:30	from cloudra.benchmark.schema import TaskSpec
cloudra/benchmark/backend.py:31	from cloudra.runtime_v2.schemas import AgentState
cloudra/benchmark/backend.py:32	
cloudra/benchmark/backend.py:33	logger = logging.getLogger(__name__)
cloudra/benchmark/backend.py:34	
cloudra/benchmark/backend.py:35	TERMINAL_STATES = frozenset({AgentState.SUCCEEDED, AgentState.FAILED, AgentState.CANCELLED})
cloudra/benchmark/backend.py:36	
cloudra/benchmark/backend.py:37	# waiting_input / waiting_approval 不是终态：benchmark 无人值守，挂起超过
cloudra/benchmark/backend.py:38	# stuck_grace 就取消该题（yolo 已开，正常路径不该出现审批等待）。
cloudra/benchmark/backend.py:39	_STUCK_STATE_VALUES = frozenset({AgentState.WAITING_INPUT.value, AgentState.WAITING_APPROVAL.value})
cloudra/benchmark/backend.py:40	
cloudra/benchmark/backend.py:41	# 最终总结里的能垒抓取：query 要求 agent 以 kcal/mol 报告能垒，取文本中
cloudra/benchmark/backend.py:42	# 最后一个 "<number> kcal/mol" 出现（总结陈述通常在最末）。
cloudra/benchmark/backend.py:43	_BARRIER_RE = re.compile(r"(-?\d+(?:\.\d+)?)\s*kcal\s*/?\s*mol", re.IGNORECASE)
cloudra/benchmark/backend.py:44	
cloudra/benchmark/backend.py:45	
cloudra/benchmark/backend.py:46	def extract_reported_barrier(text: str) -> Optional[float]:
cloudra/benchmark/backend.py:47	    """Last "<number> kcal/mol" in the agent's final summary, or None."""
cloudra/benchmark/backend.py:48	    matches = _BARRIER_RE.findall(text or "")
cloudra/benchmark/backend.py:49	    return float(matches[-1]) if matches else None
cloudra/benchmark/backend.py:50	
cloudra/benchmark/backend.py:51	
cloudra/benchmark/backend.py:52	class RuntimeV2Backend:
cloudra/benchmark/backend.py:53	    """`AgentBackend` implementation: run each task on the real runtime_v2.
cloudra/benchmark/backend.py:54	
cloudra/benchmark/backend.py:55	    `llm` / `tools` / `run_summarizer` are injectable so tests can run the
cloudra/benchmark/backend.py:56	    whole wiring hermetically; production defaults are the legacy bridges.
cloudra/benchmark/backend.py:57	    """
cloudra/benchmark/backend.py:58	
cloudra/benchmark/backend.py:59	    def __init__(
cloudra/benchmark/backend.py:60	        self,
cloudra/benchmark/backend.py:61	        *,
cloudra/benchmark/backend.py:62	        run_timeout_s: float = 7200.0,
cloudra/benchmark/backend.py:63	        stuck_grace_s: float = 300.0,
cloudra/benchmark/backend.py:64	        poll_interval_s: float = 2.0,
cloudra/benchmark/backend.py:65	        llm: Any = None,
cloudra/benchmark/backend.py:66	        tools: Any = None,
cloudra/benchmark/backend.py:67	        run_summarizer: Optional[Callable[[list[dict]], Awaitable[str]]] = None,
cloudra/benchmark/backend.py:68	        start_scheduler: bool = True,
cloudra/benchmark/backend.py:69	    ) -> None:
cloudra/benchmark/backend.py:70	        self.run_timeout_s = run_timeout_s
cloudra/benchmark/backend.py:71	        self.stuck_grace_s = stuck_grace_s
cloudra/benchmark/backend.py:72	        self.poll_interval_s = poll_interval_s
cloudra/benchmark/backend.py:73	        self.llm = llm
cloudra/benchmark/backend.py:74	        self.tools = tools
cloudra/benchmark/backend.py:75	        self.run_summarizer = run_summarizer
cloudra/benchmark/backend.py:76	        self.start_scheduler = start_scheduler
cloudra/benchmark/backend.py:77	        # 每次任务的 RuntimeCore 按 task_id 留存，供测试与调试内省事件日志；
cloudra/benchmark/backend.py:78	        # 并行跑题时多线程并发写，用锁保护。
cloudra/benchmark/backend.py:79	        self._cores: dict[str, Any] = {}
cloudra/benchmark/backend.py:80	        self._cores_lock = threading.Lock()
cloudra/benchmark/backend.py:81	
\end{lstlisting}
\clearpage
\begin{lstlisting}[firstnumber=82,numbers=left,numberstyle=\tiny\ttfamily,numbersep=5pt]
cloudra/benchmark/backend.py:82	    @property
cloudra/benchmark/backend.py:83	    def last_core(self) -> Any:
cloudra/benchmark/backend.py:84	        """最近一次任务的 core（串行调试便利）；并行时用 core_by_task()。"""
cloudra/benchmark/backend.py:85	        with self._cores_lock:
cloudra/benchmark/backend.py:86	            return next(reversed(self._cores.values())) if self._cores else None
cloudra/benchmark/backend.py:87	
cloudra/benchmark/backend.py:88	    def core_by_task(self, task_id: str) -> Any:
cloudra/benchmark/backend.py:89	        with self._cores_lock:
cloudra/benchmark/backend.py:90	            return self._cores.get(task_id)
cloudra/benchmark/backend.py:91	
cloudra/benchmark/backend.py:92	    def __call__(self, task: TaskSpec, sandbox: Path) -> Optional[float]:
cloudra/benchmark/backend.py:93	        return asyncio.run(self._run_task(task, Path(sandbox)))
cloudra/benchmark/backend.py:94	
cloudra/benchmark/backend.py:95	    async def _run_task(self, task: TaskSpec, sandbox: Path) -> Optional[float]:
cloudra/benchmark/backend.py:96	        from cloudra.runtime_v2.adapters import (
cloudra/benchmark/backend.py:97	            LegacyLLMProvider,
cloudra/benchmark/backend.py:98	            LegacyPromptProvider,
cloudra/benchmark/backend.py:99	            LegacyToolExecutor,
cloudra/benchmark/backend.py:100	            legacy_default_runtime_settings,
cloudra/benchmark/backend.py:101	        )
cloudra/benchmark/backend.py:102	        from cloudra.runtime_v2.agent_runner import AgentRunner
cloudra/benchmark/backend.py:103	        from cloudra.runtime_v2.commands import CommandBus
cloudra/benchmark/backend.py:104	        from cloudra.runtime_v2.core import RuntimeCore
cloudra/benchmark/backend.py:105	        from cloudra.runtime_v2.policies import default_stages
cloudra/benchmark/backend.py:106	        from cloudra.runtime_v2.schemas import Command
cloudra/benchmark/backend.py:107	        from cloudra.runtime_v2.turn_pipeline import TurnPipeline
cloudra/benchmark/backend.py:108	
cloudra/benchmark/backend.py:109	        sandbox.mkdir(parents=True, exist_ok=True)
cloudra/benchmark/backend.py:110	        sandbox = sandbox.resolve()
cloudra/benchmark/backend.py:111	
cloudra/benchmark/backend.py:112	        core = RuntimeCore(strict_events=True)
cloudra/benchmark/backend.py:113	        with self._cores_lock:
cloudra/benchmark/backend.py:114	            self._cores[task.task_id] = core
cloudra/benchmark/backend.py:115	        core.runtime_defaults = legacy_default_runtime_settings()
cloudra/benchmark/backend.py:116	        stage = MethodPolicyStage(task.method_policy, base_dir=sandbox)
cloudra/benchmark/backend.py:117	        pipeline = TurnPipeline(
cloudra/benchmark/backend.py:118	            event_log=core.event_log,
cloudra/benchmark/backend.py:119	            stages=[*default_stages(core), stage],
cloudra/benchmark/backend.py:120	            llm=self.llm or LegacyLLMProvider(),
cloudra/benchmark/backend.py:121	            tools=self.tools or LegacyToolExecutor(core),
cloudra/benchmark/backend.py:122	            prompt_provider=LegacyPromptProvider(),
cloudra/benchmark/backend.py:123	            event_hub=core.event_hub,
cloudra/benchmark/backend.py:124	            notices=core.notices,
cloudra/benchmark/backend.py:125	        )
cloudra/benchmark/backend.py:126	        core.runner = AgentRunner(core, pipeline, run_summarizer=self.run_summarizer)
cloudra/benchmark/backend.py:127	        scheduler = None
cloudra/benchmark/backend.py:128	        if self.start_scheduler:
cloudra/benchmark/backend.py:129	            from cloudra.scheduler.daemon import SchedulerDaemon
cloudra/benchmark/backend.py:130	
cloudra/benchmark/backend.py:131	            # worker_id 含 pid+task_id：并行跑题时同进程多个实例不能共享
cloudra/benchmark/backend.py:132	            # 默认的 scheduler:<pid>，否则租约归属互相串。
cloudra/benchmark/backend.py:133	            scheduler = SchedulerDaemon(worker_id=f"scheduler:{os.getpid()}:{task.task_id}")
\end{lstlisting}
\clearpage
\begin{lstlisting}[firstnumber=134,numbers=left,numberstyle=\tiny\ttfamily,numbersep=5pt]
cloudra/benchmark/backend.py:134	            await scheduler.start()
cloudra/benchmark/backend.py:135	
cloudra/benchmark/backend.py:136	        agent = None
cloudra/benchmark/backend.py:137	        bus = CommandBus(core)
cloudra/benchmark/backend.py:138	        try:
cloudra/benchmark/backend.py:139	            project = (
cloudra/benchmark/backend.py:140	                await bus.execute(
cloudra/benchmark/backend.py:141	                    Command(
cloudra/benchmark/backend.py:142	                        type="project.create",
cloudra/benchmark/backend.py:143	                        payload={"name": f"benchmark:{task.task_id}", "workspace_dir": str(sandbox)},
cloudra/benchmark/backend.py:144	                    )
cloudra/benchmark/backend.py:145	                )
cloudra/benchmark/backend.py:146	            ).data
cloudra/benchmark/backend.py:147	            agent = (
cloudra/benchmark/backend.py:148	                await bus.execute(
cloudra/benchmark/backend.py:149	                    Command(
cloudra/benchmark/backend.py:150	                        type="agent.spawn",
cloudra/benchmark/backend.py:151	                        target={"project_id": project.id},
cloudra/benchmark/backend.py:152	                        payload={
cloudra/benchmark/backend.py:153	                            "role_id": task.prompt.role,
cloudra/benchmark/backend.py:154	                            "query": task.prompt.query,
cloudra/benchmark/backend.py:155	                            "settings": {
cloudra/benchmark/backend.py:156	                                "orchestration_mode": task.prompt.orchestration_mode.value,
cloudra/benchmark/backend.py:157	                                # 无人值守：破坏性工具审批走 yolo 自动通道
cloudra/benchmark/backend.py:158	                                "yolo": True,
cloudra/benchmark/backend.py:159	                            },
cloudra/benchmark/backend.py:160	                        },
cloudra/benchmark/backend.py:161	                    )
cloudra/benchmark/backend.py:162	                )
cloudra/benchmark/backend.py:163	            ).data
cloudra/benchmark/backend.py:164	            await self._wait_terminal(bus, agent)
cloudra/benchmark/backend.py:165	            return self._final_barrier(core, agent.id)
cloudra/benchmark/backend.py:166	        finally:
cloudra/benchmark/backend.py:167	            if agent is not None and core.runner is not None:
cloudra/benchmark/backend.py:168	                await core.runner.cancel(agent.id)
cloudra/benchmark/backend.py:169	            # benchmark core 默认无 store（不落盘），这里是契约性的关机补快照：
cloudra/benchmark/backend.py:170	            # 将来接入持久化时无需再改。
cloudra/benchmark/backend.py:171	            await bus.flush_snapshot()
cloudra/benchmark/backend.py:172	            if scheduler is not None:
cloudra/benchmark/backend.py:173	                await scheduler.close()
cloudra/benchmark/backend.py:174	
cloudra/benchmark/backend.py:175	    async def _wait_terminal(self, bus, agent) -> None:
cloudra/benchmark/backend.py:176	        from cloudra.runtime_v2.schemas import Command
cloudra/benchmark/backend.py:177	
cloudra/benchmark/backend.py:178	        loop = asyncio.get_running_loop()
cloudra/benchmark/backend.py:179	        deadline = loop.time() + self.run_timeout_s
cloudra/benchmark/backend.py:180	        stuck_since: float | None = None
cloudra/benchmark/backend.py:181	        while agent.state not in TERMINAL_STATES:
cloudra/benchmark/backend.py:182	            now = loop.time()
cloudra/benchmark/backend.py:183	            reason = ""
cloudra/benchmark/backend.py:184	            if now > deadline:
cloudra/benchmark/backend.py:185	                reason = f"benchmark run timeout ({self.run_timeout_s}s)"
\end{lstlisting}
\clearpage
\begin{lstlisting}[firstnumber=186,numbers=left,numberstyle=\tiny\ttfamily,numbersep=5pt]
cloudra/benchmark/backend.py:186	            elif agent.state.value in _STUCK_STATE_VALUES:
cloudra/benchmark/backend.py:187	                stuck_since = stuck_since if stuck_since is not None else now
cloudra/benchmark/backend.py:188	                if now - stuck_since > self.stuck_grace_s:
cloudra/benchmark/backend.py:189	                    reason = f"agent stuck in {agent.state.value} with nobody to answer"
cloudra/benchmark/backend.py:190	            else:
cloudra/benchmark/backend.py:191	                stuck_since = None
cloudra/benchmark/backend.py:192	            if reason:
cloudra/benchmark/backend.py:193	                logger.warning("cancelling benchmark agent %s: %s", agent.id, reason)
cloudra/benchmark/backend.py:194	                await bus.execute(
cloudra/benchmark/backend.py:195	                    Command(
cloudra/benchmark/backend.py:196	                        type="agent.cancel",
cloudra/benchmark/backend.py:197	                        target={"session_id": agent.id},
cloudra/benchmark/backend.py:198	                        payload={"message": f"[benchmark] cancelled: {reason}"},
cloudra/benchmark/backend.py:199	                    )
cloudra/benchmark/backend.py:200	                )
cloudra/benchmark/backend.py:201	                break
cloudra/benchmark/backend.py:202	            await asyncio.sleep(self.poll_interval_s)
cloudra/benchmark/backend.py:203	
cloudra/benchmark/backend.py:204	    def _final_barrier(self, core, session_id: str) -> Optional[float]:
cloudra/benchmark/backend.py:205	        events = core.event_log.query_session(session_id)
cloudra/benchmark/backend.py:206	        for envelope in reversed(events):
cloudra/benchmark/backend.py:207	            if envelope.type == "agent.message_emitted":
cloudra/benchmark/backend.py:208	                content = str((envelope.data or {}).get("content") or "")
cloudra/benchmark/backend.py:209	                barrier = extract_reported_barrier(content)
cloudra/benchmark/backend.py:210	                if barrier is None:
cloudra/benchmark/backend.py:211	                    logger.warning("agent %s final message reports no barrier", session_id)
cloudra/benchmark/backend.py:212	                return barrier
cloudra/benchmark/backend.py:213	        return None
cloudra/benchmark/gen_easy_tasks.py:1	"""Generate benchmark task JSONs from ideapool/reaction_candidates/easy.md.
cloudra/benchmark/gen_easy_tasks.py:2	
cloudra/benchmark/gen_easy_tasks.py:3	Reads the curated easy-tier markdown and emits one TaskSpec JSON per task into
cloudra/benchmark/gen_easy_tasks.py:4	ideapool/tasks/easy/<task_id>.json. The markdown is the human-curated source
cloudra/benchmark/gen_easy_tasks.py:5	of truth; this generator is deterministic and re-runnable.
cloudra/benchmark/gen_easy_tasks.py:6	
cloudra/benchmark/gen_easy_tasks.py:7	Usage:  python -m cloudra.benchmark.gen_easy_tasks [--out DIR]
cloudra/benchmark/gen_easy_tasks.py:8	"""
cloudra/benchmark/gen_easy_tasks.py:9	
cloudra/benchmark/gen_easy_tasks.py:10	from __future__ import annotations
cloudra/benchmark/gen_easy_tasks.py:11	
cloudra/benchmark/gen_easy_tasks.py:12	import argparse
cloudra/benchmark/gen_easy_tasks.py:13	import json
cloudra/benchmark/gen_easy_tasks.py:14	import re
cloudra/benchmark/gen_easy_tasks.py:15	from pathlib import Path
cloudra/benchmark/gen_easy_tasks.py:16	
cloudra/benchmark/gen_easy_tasks.py:17	from cloudra.benchmark.schema import (
cloudra/benchmark/gen_easy_tasks.py:18	    BarrierReference,
cloudra/benchmark/gen_easy_tasks.py:19	    BarrierType,
cloudra/benchmark/gen_easy_tasks.py:20	    BondChangeSet,
cloudra/benchmark/gen_easy_tasks.py:21	    GoldMechanismGraph,
cloudra/benchmark/gen_easy_tasks.py:22	    GoldSpec,
cloudra/benchmark/gen_easy_tasks.py:23	    MechanismEdge,
cloudra/benchmark/gen_easy_tasks.py:24	    MechanismNode,
\end{lstlisting}
\clearpage
\begin{lstlisting}[firstnumber=25,numbers=left,numberstyle=\tiny\ttfamily,numbersep=5pt]
cloudra/benchmark/gen_easy_tasks.py:25	    NodeRole,
cloudra/benchmark/gen_easy_tasks.py:26	    PromptSpec,
cloudra/benchmark/gen_easy_tasks.py:27	    ReferenceComputation,
cloudra/benchmark/gen_easy_tasks.py:28	    RubricSpec,
cloudra/benchmark/gen_easy_tasks.py:29	    SystemSpec,
cloudra/benchmark/gen_easy_tasks.py:30	    TaskSpec,
cloudra/benchmark/gen_easy_tasks.py:31	    Tier,
cloudra/benchmark/gen_easy_tasks.py:32	    VerificationStatus,
cloudra/benchmark/gen_easy_tasks.py:33	)
cloudra/benchmark/gen_easy_tasks.py:34	
cloudra/benchmark/gen_easy_tasks.py:35	REPO_ROOT = Path(__file__).resolve().parents[2]
cloudra/benchmark/gen_easy_tasks.py:36	DEFAULT_SRC = REPO_ROOT / "ideapool" / "reaction_candidates" / "easy.md"
cloudra/benchmark/gen_easy_tasks.py:37	DEFAULT_OUT = REPO_ROOT / "ideapool" / "tasks" / "easy"
cloudra/benchmark/gen_easy_tasks.py:38	
cloudra/benchmark/gen_easy_tasks.py:39	ROW_RE = re.compile(r"^\|\s*([A-Z]\d+)\s*\|")
cloudra/benchmark/gen_easy_tasks.py:40	# SMILES cell like `C=CC=C.C=C>>C1CC=CCC1`，0/1  (backticks optional).
cloudra/benchmark/gen_easy_tasks.py:41	# charge may use the Unicode minus '−'; SMILES may carry a trailing note
cloudra/benchmark/gen_easy_tasks.py:42	# like `O=NO>>O=NO`（cis/trans），0/1 before the q/mult separator.
cloudra/benchmark/gen_easy_tasks.py:43	SMILES_RE = re.compile(r"`([^`]+)`(?:（[^）]*）)?\s*[，,]\s*([+−\-]?\d+)\s*/\s*(\d+)")
cloudra/benchmark/gen_easy_tasks.py:44	VERIFY_RE = re.compile(r"\[(已核实存在|已核实|部分核实|未核实)\]")
cloudra/benchmark/gen_easy_tasks.py:45	ELEMENTS_RE = re.compile(r"^([A-Z][a-z]?(?:,[A-Z][a-z]?)*)")
cloudra/benchmark/gen_easy_tasks.py:46	HEAVY_RE = re.compile(r"(\d+)\s*重原子")
cloudra/benchmark/gen_easy_tasks.py:47	
cloudra/benchmark/gen_easy_tasks.py:48	_VERIFY_MAP = {
cloudra/benchmark/gen_easy_tasks.py:49	    "已核实": VerificationStatus.VERIFIED,
cloudra/benchmark/gen_easy_tasks.py:50	    "已核实存在": VerificationStatus.VERIFIED,
cloudra/benchmark/gen_easy_tasks.py:51	    "部分核实": VerificationStatus.PARTIAL,
cloudra/benchmark/gen_easy_tasks.py:52	    "未核实": VerificationStatus.UNVERIFIED,
cloudra/benchmark/gen_easy_tasks.py:53	}
cloudra/benchmark/gen_easy_tasks.py:54	
cloudra/benchmark/gen_easy_tasks.py:55	_BARRIER_TYPE_MAP = {
cloudra/benchmark/gen_easy_tasks.py:56	    "ΔH‡": BarrierType.DH,
cloudra/benchmark/gen_easy_tasks.py:57	    "ΔE‡": BarrierType.DE,
cloudra/benchmark/gen_easy_tasks.py:58	    "ΔG‡": BarrierType.DG,
cloudra/benchmark/gen_easy_tasks.py:59	    "Ea": BarrierType.EA,
cloudra/benchmark/gen_easy_tasks.py:60	}
cloudra/benchmark/gen_easy_tasks.py:61	
cloudra/benchmark/gen_easy_tasks.py:62	# 能垒口径（barrier_reference）统一为 CENTRAL（central_complex：TS 相对反应物
cloudra/benchmark/gen_easy_tasks.py:63	# 络合物/预反应络合物——双分子离子-分子反应即离子-偶极络合物；单分子与分子内
cloudra/benchmark/gen_easy_tasks.py:64	# 过程的反应物阱本身就是"反应物络合物"）。easy 档不采用相对分离反应物渐近线
cloudra/benchmark/gen_easy_tasks.py:65	# 的口径；个别题若日后确需 reactants 口径，在此按 task_id 显式登记覆盖。能垒
cloudra/benchmark/gen_easy_tasks.py:66	# 缺测时 barrier_reference 记 None。
cloudra/benchmark/gen_easy_tasks.py:67	_BARRIER_REFERENCE: dict[str, BarrierReference] = {}
cloudra/benchmark/gen_easy_tasks.py:68	
cloudra/benchmark/gen_easy_tasks.py:69	# 口径核对登记：下列题的文献锚点明确是相对分离反应物渐近线的数值，且无络合物
cloudra/benchmark/gen_easy_tasks.py:70	# 口径的文献值可替换——数值原样保留（中性预反应络合物的 vdW 阱深一般只有
cloudra/benchmark/gen_easy_tasks.py:71	# ~1-3 kcal/mol，锚点在 accuracy 层高斯容差内仍有参考意义），verification
cloudra/benchmark/gen_easy_tasks.py:72	# 一律降为 unverified 待核，绝不按新口径编造数值。其余题的反应物阱即络合物
cloudra/benchmark/gen_easy_tasks.py:73	# 口径（单分子/分子内题）或出处本来就是中央能垒（C1/E1），直接挂 CENTRAL。
cloudra/benchmark/gen_easy_tasks.py:74	_BARRIER_FRAME_UNVERIFIED = {
cloudra/benchmark/gen_easy_tasks.py:75	    "D1", "D3", "D4",   # 双分子 DA：实验 Ea / 理论 ΔE‡ 均为分离反应物口径
cloudra/benchmark/gen_easy_tasks.py:76	    "C2", "C3", "F1",   # 题源单元格明确写"TS 相对反应物渐近线"
\end{lstlisting}
\clearpage
\begin{lstlisting}[firstnumber=77,numbers=left,numberstyle=\tiny\ttfamily,numbersep=5pt]
cloudra/benchmark/gen_easy_tasks.py:77	    "F3", "T1",         # 微溶剂化：文献能垒相对分离单体（"能垒降 ~30%"的口径）
cloudra/benchmark/gen_easy_tasks.py:78	    "G1", "G2", "G3",   # 双分子协同：文献值相对分离反应物（G3 题源本就未核实）
cloudra/benchmark/gen_easy_tasks.py:79	}
cloudra/benchmark/gen_easy_tasks.py:80	
cloudra/benchmark/gen_easy_tasks.py:81	# 条件列中译英（按原字符串映射；新增条件必须在此登记，否则生成直接报错）。
cloudra/benchmark/gen_easy_tasks.py:82	# M 组表第 4 列是过程类型而非条件，Easy 档统一定义为气相可算，故 M 题
cloudra/benchmark/gen_easy_tasks.py:83	# conditions 一律 "gas phase"，过程类型并入 step_type 英译。
cloudra/benchmark/gen_easy_tasks.py:84	_CONDITIONS_EN = {
cloudra/benchmark/gen_easy_tasks.py:85	    "气相": "gas phase",
cloudra/benchmark/gen_easy_tasks.py:86	    "气相 487–648 °C 动力学": "gas phase, 487-648 degC kinetics",
cloudra/benchmark/gen_easy_tasks.py:87	    "气相/液相室温慢": "gas/liquid phase, slow at room temperature",
cloudra/benchmark/gen_easy_tasks.py:88	    "气相/烃溶剂加热": "gas phase or hydrocarbon solvent, heated",
cloudra/benchmark/gen_easy_tasks.py:89	    "气相 ~250 °C": "gas phase, ~250 degC",
cloudra/benchmark/gen_easy_tasks.py:90	    "气相/溶液室温即快": "gas phase or solution, fast at room temperature",
cloudra/benchmark/gen_easy_tasks.py:91	    "气相 ~200 °C": "gas phase, ~200 degC",
cloudra/benchmark/gen_easy_tasks.py:92	    "气相 403–448 K": "gas phase, 403-448 K",
cloudra/benchmark/gen_easy_tasks.py:93	    "气相 390–434 K": "gas phase, 390-434 K",
cloudra/benchmark/gen_easy_tasks.py:94	    "气相 ~350 °C": "gas phase, ~350 degC",
cloudra/benchmark/gen_easy_tasks.py:95	    "气相/溶液": "gas phase or solution",
cloudra/benchmark/gen_easy_tasks.py:96	    "气相/水溶液": "gas phase or aqueous solution",
cloudra/benchmark/gen_easy_tasks.py:97	    "气相微溶剂化": "gas phase, microsolvated",
cloudra/benchmark/gen_easy_tasks.py:98	    "气相/星际": "gas phase / interstellar",
cloudra/benchmark/gen_easy_tasks.py:99	    "气相高温": "gas phase, high temperature",
cloudra/benchmark/gen_easy_tasks.py:100	}
cloudra/benchmark/gen_easy_tasks.py:101	
cloudra/benchmark/gen_easy_tasks.py:102	# 题名英译，按 task_id 映射；新增题目必须登记（漏译 KeyError 直接报错）。
cloudra/benchmark/gen_easy_tasks.py:103	_NAME_EN = {
cloudra/benchmark/gen_easy_tasks.py:104	    "D1": "butadiene + ethylene Diels-Alder (parent DA)",
cloudra/benchmark/gen_easy_tasks.py:105	    "D3": "cyclopentadiene dimerization (endo, self DA)",
cloudra/benchmark/gen_easy_tasks.py:106	    "D4": "cyclopentadiene + acrolein DA (heteroatom dienophile)",
cloudra/benchmark/gen_easy_tasks.py:107	    "P1": "1,5-hexadiene Cope rearrangement ([3,3], degenerate)",
cloudra/benchmark/gen_easy_tasks.py:108	    "P2": "cyclopentadiene [1,5]-H shift (degenerate)",
cloudra/benchmark/gen_easy_tasks.py:109	    "P3": "allyl vinyl ether Claisen rearrangement ([3,3])",
cloudra/benchmark/gen_easy_tasks.py:110	    "L1": "cyclobutene -> butadiene conrotatory ring opening",
cloudra/benchmark/gen_easy_tasks.py:111	    "L2": "hexatriene -> cyclohexadiene disrotatory ring closure",
cloudra/benchmark/gen_easy_tasks.py:112	    "X1": "beta-propiolactone -> ethylene + CO2 (decarboxylative extrusion)",
cloudra/benchmark/gen_easy_tasks.py:113	    "C1": "F- + CH3Cl -> CH3F + Cl-",
cloudra/benchmark/gen_easy_tasks.py:114	    "C2": "CN- + CH3Cl -> CH3CN + Cl-",
cloudra/benchmark/gen_easy_tasks.py:115	    "C3": "F- + CH3Cl double inversion (retention pathway)",
cloudra/benchmark/gen_easy_tasks.py:116	    "E1": "F- + CH3CH2Cl anti-E2 -> ethylene + HF + Cl-",
cloudra/benchmark/gen_easy_tasks.py:117	    "F1": "BH4- + formaldehyde -> methoxide + BH3",
cloudra/benchmark/gen_easy_tasks.py:118	    "F2": "CN- + formaldehyde -> cyanohydrin alkoxide",
cloudra/benchmark/gen_easy_tasks.py:119	    "F3": "formaldehyde + 2 H2O -> methanediol + H2O (water-assisted hydration)",
cloudra/benchmark/gen_easy_tasks.py:120	    "T1": "acetaldehyde + H2O -> vinyl alcohol + H2O (water-assisted tautomerization)",
cloudra/benchmark/gen_easy_tasks.py:121	    "T3": "HNC -> HCN (1,2-H shift)",
cloudra/benchmark/gen_easy_tasks.py:122	    "G1": "BH3 + ethylene hydroboration (concerted syn)",
cloudra/benchmark/gen_easy_tasks.py:123	    "G2": "diimide N2H2 cis reduction of ethylene (double H transfer)",
cloudra/benchmark/gen_easy_tasks.py:124	    "G3": "propylene + ethylene ene reaction (Alder-ene)",
cloudra/benchmark/gen_easy_tasks.py:125	    "M1": "2,2'-dimethylbiphenyl atropisomer rotation",
cloudra/benchmark/gen_easy_tasks.py:126	    "M3": "NH3 umbrella inversion",
cloudra/benchmark/gen_easy_tasks.py:127	    "M6": "ethane rotation (staggered<->eclipsed)",
cloudra/benchmark/gen_easy_tasks.py:128	    "M8": "cyclohexane chair flip",
\end{lstlisting}
\clearpage
\begin{lstlisting}[firstnumber=129,numbers=left,numberstyle=\tiny\ttfamily,numbersep=5pt]
cloudra/benchmark/gen_easy_tasks.py:129	    "M10": "acetaldehyde methyl rotation",
cloudra/benchmark/gen_easy_tasks.py:130	    "M12": "H5O2+ Zundel proton transfer",
cloudra/benchmark/gen_easy_tasks.py:131	    "M14": "acrolein s-trans <-> s-cis rotation",
cloudra/benchmark/gen_easy_tasks.py:132	    "M15": "HONO cis-trans isomerization",
cloudra/benchmark/gen_easy_tasks.py:133	    "M16": "formic acid dehydrogenation HCOOH -> CO2 + H2",
cloudra/benchmark/gen_easy_tasks.py:134	}
cloudra/benchmark/gen_easy_tasks.py:135	
cloudra/benchmark/gen_easy_tasks.py:136	# gold 机理概要（step_type）英译，按 task_id 映射；新增题目必须登记。
cloudra/benchmark/gen_easy_tasks.py:137	_STEP_TYPE_EN = {
cloudra/benchmark/gen_easy_tasks.py:138	    "D1": "concerted [4+2] cycloaddition, single step via symmetric TS",
cloudra/benchmark/gen_easy_tasks.py:139	    "D3": "concerted [4+2], bispericyclic TS; one molecule as diene, one as dienophile",
cloudra/benchmark/gen_easy_tasks.py:140	    "D4": "concerted [4+2], endo/exo selectivity, aldehyde dienophile",
cloudra/benchmark/gen_easy_tasks.py:141	    "P1": "[3,3]-sigmatropic shift, chair TS, degenerate rearrangement",
cloudra/benchmark/gen_easy_tasks.py:142	    "P2": "[1,5]-sigmatropic H shift, suprafacially allowed, unimolecular",
cloudra/benchmark/gen_easy_tasks.py:143	    "P3": "[3,3]-sigmatropic shift, chair TS, Claisen (O-containing complement to all-carbon Cope P1) -> pent-4-enal",
cloudra/benchmark/gen_easy_tasks.py:144	    "L1": "4pi electrocyclic ring opening, conrotatory, single step",
cloudra/benchmark/gen_easy_tasks.py:145	    "L2": "6pi electrocyclic ring closure, disrotatory, single step",
cloudra/benchmark/gen_easy_tasks.py:146	    "X1": "concerted asynchronous C-C/C-O cleavage, CO2 extrusion",
cloudra/benchmark/gen_easy_tasks.py:147	    "C1": "backside attack, single step, exothermic",
cloudra/benchmark/gen_easy_tasks.py:148	    "C2": "backside attack, single step (C-end attack of ambident nucleophile)",
cloudra/benchmark/gen_easy_tasks.py:149	    "C3": "double inversion (retention): proton-abstraction-induced inversion to a [FH...CH2Cl]- intermediate, then a second inversion; gold models the rate-determining Cs-symmetric TS as a single step",
cloudra/benchmark/gen_easy_tasks.py:150	    "E1": "anti-periplanar concerted elimination, single step (competes with SN2)",
cloudra/benchmark/gen_easy_tasks.py:151	    "F1": "hydride transfer, single step via asynchronous TS",
cloudra/benchmark/gen_easy_tasks.py:152	    "F2": "anion nucleophilic addition to carbonyl, single step (Burgi-Dunitz trajectory ~107 deg)",
cloudra/benchmark/gen_easy_tasks.py:153	    "F3": "extra water as proton shuttle, six-center cyclic TS, single step",
cloudra/benchmark/gen_easy_tasks.py:154	    "T1": "water as proton shuttle, six-center cyclic TS, single step",
cloudra/benchmark/gen_easy_tasks.py:155	    "T3": "intramolecular 1,2-H shift; anomalous IRC - separate C-H cleavage and N-H formation transition regions with an all-real-frequency quasi-state between",
cloudra/benchmark/gen_easy_tasks.py:156	    "G1": "four-center cyclic TS, concerted C-B and C-H formation with B-H cleavage, syn addition; nearly barrierless in gas phase, TST failure (kinetic anomaly)",
cloudra/benchmark/gen_easy_tasks.py:157	    "G2": "concerted [sigma2s+sigma2s+pi2s] double H transfer, syn stereospecific, N2 release",
cloudra/benchmark/gen_easy_tasks.py:158	    "G3": "concerted [sigma2s+pi2s+pi2s], allylic H shift + C-C formation + pi shift",
cloudra/benchmark/gen_easy_tasks.py:159	    "M1": "atropisomerism: ortho-methyl hindered rotation, Ra<->Sa via near-planar TS",
cloudra/benchmark/gen_easy_tasks.py:160	    "M3": "N inversion: double-well inversion via planar D3h TS, tunneling splitting",
cloudra/benchmark/gen_easy_tasks.py:161	    "M6": "conformational rotation: C-C torsion, threefold-symmetric TS",
cloudra/benchmark/gen_easy_tasks.py:162	    "M8": "conformational inversion: chair-chair interconversion via half-chair TS",
cloudra/benchmark/gen_easy_tasks.py:163	    "M10": "methyl rotation: threefold methyl torsion, focal-point benchmark",
cloudra/benchmark/gen_easy_tasks.py:164	    "M12": "proton cluster transfer: shared proton between two waters; O-O distance decides single/double well",
cloudra/benchmark/gen_easy_tasks.py:165	    "M14": "conformational rotation: C(sp2)-C(sp2) single-bond torsion, trans more stable",
cloudra/benchmark/gen_easy_tasks.py:166	    "M15": "conformational isomerization: N-O torsion, trans ground state, cis ~0.5-0.6 kcal/mol higher",
cloudra/benchmark/gen_easy_tasks.py:167	    "M16": "dehydrogenation: concerted H2 elimination to CO2 (competes with dehydration to CO+H2O)",
cloudra/benchmark/gen_easy_tasks.py:168	}
cloudra/benchmark/gen_easy_tasks.py:169	
cloudra/benchmark/gen_easy_tasks.py:170	# ---------------------------------------------------------------------------
cloudra/benchmark/gen_easy_tasks.py:171	# 原子映射反应物 SMILES（query 专用）与 gold 边键变化，按 task_id 登记；
cloudra/benchmark/gen_easy_tasks.py:172	# 新增题目两处都必须登记（漏登 KeyError 直接报错）。
cloudra/benchmark/gen_easy_tasks.py:173	#
cloudra/benchmark/gen_easy_tasks.py:174	# 索引约定：query 向 agent 展示 _MAPPED_REACTANTS 的原子映射反应物 SMILES
cloudra/benchmark/gen_easy_tasks.py:175	# （RDKit 可读的 1-based atom-map 编号；参与键变化的 H 以显式 [H:n] 给出，
cloudra/benchmark/gen_easy_tasks.py:176	# 其余 H 保持隐式计数），gold 与 agent 交付边的 bond_changes [i, j] 对都引用
cloudra/benchmark/gen_easy_tasks.py:177	# 这套映射号，两侧因此共用同一套原子编号。system.reactants/products 与 gold
cloudra/benchmark/gen_easy_tasks.py:178	# 节点 SMILES 保持无映射（工具链与完整性取证用无映射形式）。
cloudra/benchmark/gen_easy_tasks.py:179	# 键变化口径：break = 该基元步内键级下降（含断至 0），form = 键级上升（含自
cloudra/benchmark/gen_easy_tasks.py:180	# 0 新建）；周环反应的 π 迁移据此记为旧 π 键 break + 新 π 键 form，双键变
\end{lstlisting}
\clearpage
\begin{lstlisting}[firstnumber=181,numbers=left,numberstyle=\tiny\ttfamily,numbersep=5pt]
cloudra/benchmark/gen_easy_tasks.py:181	# 单键（如羰基受亲核进攻）记 break。纯构象/翻转步（M1/M3/M6/M8/M10/M14/
cloudra/benchmark/gen_easy_tasks.py:182	# M15，扭转或伞形翻转，无键级变化）两个列表皆空。
cloudra/benchmark/gen_easy_tasks.py:183	# 注意：RDKit 默认 MolFromSmiles 会移除显式 H，核对映射 H 需 removeHs=False。
cloudra/benchmark/gen_easy_tasks.py:184	# ---------------------------------------------------------------------------
cloudra/benchmark/gen_easy_tasks.py:185	_MAPPED_REACTANTS = {
cloudra/benchmark/gen_easy_tasks.py:186	    "D1": "[CH2:1]=[CH:2][CH:3]=[CH2:4].[CH2:5]=[CH2:6]",
cloudra/benchmark/gen_easy_tasks.py:187	    "D3": "[CH:1]1=[CH:2][CH2:3][CH:4]=[CH:5]1.[CH:6]1=[CH:7][CH2:8][CH:9]=[CH:10]1",
cloudra/benchmark/gen_easy_tasks.py:188	    "D4": "[CH:1]1=[CH:2][CH2:3][CH:4]=[CH:5]1.[CH2:6]=[CH:7][CH:8]=[O:9]",
cloudra/benchmark/gen_easy_tasks.py:189	    "P1": "[CH2:1]=[CH:2][CH2:3][CH2:4][CH:5]=[CH2:6]",
cloudra/benchmark/gen_easy_tasks.py:190	    "P2": "[CH:1]1=[CH:2][CH:3]([H:6])[CH:4]=[CH:5]1",
cloudra/benchmark/gen_easy_tasks.py:191	    "P3": "[CH2:1]=[CH:2][CH2:3][O:4][CH:5]=[CH2:6]",
cloudra/benchmark/gen_easy_tasks.py:192	    "L1": "[CH:1]1=[CH:2][CH2:3][CH2:4]1",
cloudra/benchmark/gen_easy_tasks.py:193	    "L2": "[CH2:1]=[CH:2][CH:3]=[CH:4][CH:5]=[CH2:6]",
cloudra/benchmark/gen_easy_tasks.py:194	    "X1": "[O:1]=[C:2]1[O:3][CH2:4][CH2:5]1",
cloudra/benchmark/gen_easy_tasks.py:195	    "C1": "[F-:1].[CH3:2][Cl:3]",
cloudra/benchmark/gen_easy_tasks.py:196	    "C2": "[C-:1]#[N:2].[CH3:3][Cl:4]",
cloudra/benchmark/gen_easy_tasks.py:197	    "C3": "[F-:1].[CH3:2][Cl:3]",
cloudra/benchmark/gen_easy_tasks.py:198	    "E1": "[F-:1].[CH2:2]([H:4])[CH2:3][Cl:5]",
cloudra/benchmark/gen_easy_tasks.py:199	    "F1": "[BH3-:1][H:2].[CH2:3]=[O:4]",
cloudra/benchmark/gen_easy_tasks.py:200	    "F2": "[C-:1]#[N:2].[CH2:3]=[O:4]",
cloudra/benchmark/gen_easy_tasks.py:201	    "F3": "[CH2:1]=[O:2].[OH:3]([H:4]).[OH:5]([H:6])",
cloudra/benchmark/gen_easy_tasks.py:202	    "T1": "[CH2:1]([H:5])[CH:2]=[O:3].[OH:4]([H:6])",
cloudra/benchmark/gen_easy_tasks.py:203	    "T3": "[C-:1]#[N+:2][H:3]",
cloudra/benchmark/gen_easy_tasks.py:204	    "G1": "[BH2:1][H:3].[CH2:4]=[CH2:5]",
cloudra/benchmark/gen_easy_tasks.py:205	    "G2": "[N:1]([H:3])=[N:2]([H:4]).[CH2:5]=[CH2:6]",
cloudra/benchmark/gen_easy_tasks.py:206	    "G3": "[CH2:1]=[CH:2][CH2:3]([H:6]).[CH2:4]=[CH2:5]",
cloudra/benchmark/gen_easy_tasks.py:207	    "M1": "[CH3:1][c:2]1[cH:3][cH:4][cH:5][cH:6][c:7]1[c:8]1[c:9]([CH3:10])[cH:11][cH:12][cH:13][cH:14]1",
cloudra/benchmark/gen_easy_tasks.py:208	    "M3": "[NH3:1]",
cloudra/benchmark/gen_easy_tasks.py:209	    "M6": "[CH3:1][CH3:2]",
cloudra/benchmark/gen_easy_tasks.py:210	    "M8": "[CH2:1]1[CH2:2][CH2:3][CH2:4][CH2:5][CH2:6]1",
cloudra/benchmark/gen_easy_tasks.py:211	    "M10": "[CH3:1][CH:2]=[O:3]",
cloudra/benchmark/gen_easy_tasks.py:212	    "M12": "[OH2+:1][H:3].[OH2:2]",
cloudra/benchmark/gen_easy_tasks.py:213	    "M14": "[CH2:1]=[CH:2][CH:3]=[O:4]",
cloudra/benchmark/gen_easy_tasks.py:214	    "M15": "[O:1]=[N:2][OH:3]",
cloudra/benchmark/gen_easy_tasks.py:215	    "M16": "[O:1]([H:4])[C:2]([H:5])=[O:3]",
cloudra/benchmark/gen_easy_tasks.py:216	}
cloudra/benchmark/gen_easy_tasks.py:217	
cloudra/benchmark/gen_easy_tasks.py:218	# (break, form)；每个 [i, j] 对引用上方同一 task 的映射号。逐题化学依据：
cloudra/benchmark/gen_easy_tasks.py:219	# - D/P/L/X 周环：断旧 π/σ、成新 π/σ（DA 成两个新 σ + 一个 π；电环化开/闭环
cloudra/benchmark/gen_easy_tasks.py:220	#   断/成一个 σ 并整体移 π；X1 断 C-C 与 C-O、成 C=O 第二键与 C=C）。
cloudra/benchmark/gen_easy_tasks.py:221	# - C1/C3 SN2（含双反转净结果）：断 C-Cl、成 C-F；C2 断 C-Cl、成 C-C（C 端）。
cloudra/benchmark/gen_easy_tasks.py:222	# - E1：断 β-C-H 与 C-Cl，成 H-F 与 C=C π。
cloudra/benchmark/gen_easy_tasks.py:223	# - F1 断 B-H 与 C=O π、成 C-H；F2 成 C-C、断 C=O π；F3 六元环 TS：C=O π 断、
cloudra/benchmark/gen_easy_tasks.py:224	#   两个 O-H 断/成各一次（水梭）、成新 C-O。
cloudra/benchmark/gen_easy_tasks.py:225	# - T1 水梭互变：断 α-C-H 与水 O-H，成两个 O-H 与 C=C π、断 C=O π；T3 断
cloudra/benchmark/gen_easy_tasks.py:226	#   N-H、成 C-H；M12 断 O-H、成另一 O-H（Zundel 质子转移）。
cloudra/benchmark/gen_easy_tasks.py:227	# - G1 断 B-H 与 C=C π，成 C-B 与 C-H；G2 双 H 转移：断两个 N-H 与 C=C π，
cloudra/benchmark/gen_easy_tasks.py:228	#   成两个 C-H 与 N≡N 第三键；G3 ene：断烯丙基 C-H、ene π 与 enophile π，
cloudra/benchmark/gen_easy_tasks.py:229	#   成新 C-H、新 C-C σ 与迁移 π。
cloudra/benchmark/gen_easy_tasks.py:230	# - M16 脱氢：断 O-H 与 C-H，成 H-H 与 C-O 第二键（CO2）。
cloudra/benchmark/gen_easy_tasks.py:231	# - M1/M3/M6/M8/M10/M14/M15：扭转/翻转无键级变化，([] , [])。
cloudra/benchmark/gen_easy_tasks.py:232	_BOND_CHANGES: dict[str, tuple[list[list[int]], list[list[int]]]] = {
\end{lstlisting}
\clearpage
\begin{lstlisting}[firstnumber=233,numbers=left,numberstyle=\tiny\ttfamily,numbersep=5pt]
cloudra/benchmark/gen_easy_tasks.py:233	    "D1": ([[1, 2], [3, 4], [5, 6]], [[1, 6], [2, 3], [4, 5]]),
cloudra/benchmark/gen_easy_tasks.py:234	    "D3": ([[1, 2], [4, 5], [6, 7]], [[1, 5], [2, 6], [4, 7]]),
cloudra/benchmark/gen_easy_tasks.py:235	    "D4": ([[1, 2], [4, 5], [6, 7]], [[1, 5], [2, 6], [4, 7]]),
cloudra/benchmark/gen_easy_tasks.py:236	    "P1": ([[1, 2], [3, 4], [5, 6]], [[1, 6], [2, 3], [4, 5]]),
cloudra/benchmark/gen_easy_tasks.py:237	    "P2": ([[1, 2], [3, 6], [4, 5]], [[1, 5], [2, 3], [4, 6]]),
cloudra/benchmark/gen_easy_tasks.py:238	    "P3": ([[1, 2], [3, 4], [5, 6]], [[1, 6], [2, 3], [4, 5]]),
cloudra/benchmark/gen_easy_tasks.py:239	    "L1": ([[1, 2], [3, 4]], [[1, 4], [2, 3]]),
cloudra/benchmark/gen_easy_tasks.py:240	    "L2": ([[1, 2], [3, 4], [5, 6]], [[1, 6], [2, 3], [4, 5]]),
cloudra/benchmark/gen_easy_tasks.py:241	    "X1": ([[2, 5], [3, 4]], [[2, 3], [4, 5]]),
cloudra/benchmark/gen_easy_tasks.py:242	    "C1": ([[2, 3]], [[1, 2]]),
cloudra/benchmark/gen_easy_tasks.py:243	    "C2": ([[3, 4]], [[1, 3]]),
cloudra/benchmark/gen_easy_tasks.py:244	    "C3": ([[2, 3]], [[1, 2]]),
cloudra/benchmark/gen_easy_tasks.py:245	    "E1": ([[2, 4], [3, 5]], [[1, 4], [2, 3]]),
cloudra/benchmark/gen_easy_tasks.py:246	    "F1": ([[1, 2], [3, 4]], [[2, 3]]),
cloudra/benchmark/gen_easy_tasks.py:247	    "F2": ([[3, 4]], [[1, 3]]),
cloudra/benchmark/gen_easy_tasks.py:248	    "F3": ([[1, 2], [3, 4], [5, 6]], [[1, 3], [2, 6], [4, 5]]),
cloudra/benchmark/gen_easy_tasks.py:249	    "T1": ([[1, 5], [2, 3], [4, 6]], [[1, 2], [3, 6], [4, 5]]),
cloudra/benchmark/gen_easy_tasks.py:250	    "T3": ([[2, 3]], [[1, 3]]),
cloudra/benchmark/gen_easy_tasks.py:251	    "G1": ([[1, 3], [4, 5]], [[1, 4], [3, 5]]),
cloudra/benchmark/gen_easy_tasks.py:252	    "G2": ([[1, 3], [2, 4], [5, 6]], [[1, 2], [3, 5], [4, 6]]),
cloudra/benchmark/gen_easy_tasks.py:253	    "G3": ([[1, 2], [3, 6], [4, 5]], [[1, 5], [2, 3], [4, 6]]),
cloudra/benchmark/gen_easy_tasks.py:254	    "M1": ([], []),
cloudra/benchmark/gen_easy_tasks.py:255	    "M3": ([], []),
cloudra/benchmark/gen_easy_tasks.py:256	    "M6": ([], []),
cloudra/benchmark/gen_easy_tasks.py:257	    "M8": ([], []),
cloudra/benchmark/gen_easy_tasks.py:258	    "M10": ([], []),
cloudra/benchmark/gen_easy_tasks.py:259	    "M12": ([[1, 3]], [[2, 3]]),
cloudra/benchmark/gen_easy_tasks.py:260	    "M14": ([], []),
cloudra/benchmark/gen_easy_tasks.py:261	    "M15": ([], []),
cloudra/benchmark/gen_easy_tasks.py:262	    "M16": ([[1, 4], [2, 5]], [[1, 2], [4, 5]]),
cloudra/benchmark/gen_easy_tasks.py:263	}
cloudra/benchmark/gen_easy_tasks.py:264	
cloudra/benchmark/gen_easy_tasks.py:265	
cloudra/benchmark/gen_easy_tasks.py:266	def _split_cells(line: str) -> list[str]:
cloudra/benchmark/gen_easy_tasks.py:267	    cells = [c.strip() for c in line.strip().strip("|").split("|")]
cloudra/benchmark/gen_easy_tasks.py:268	    return cells
cloudra/benchmark/gen_easy_tasks.py:269	
cloudra/benchmark/gen_easy_tasks.py:270	
cloudra/benchmark/gen_easy_tasks.py:271	def _parse_smiles_cell(cell: str) -> tuple[list[str], list[str], int, int] | None:
cloudra/benchmark/gen_easy_tasks.py:272	    m = SMILES_RE.search(cell)
cloudra/benchmark/gen_easy_tasks.py:273	    if not m:
cloudra/benchmark/gen_easy_tasks.py:274	        return None
cloudra/benchmark/gen_easy_tasks.py:275	    rxn, q_raw, mult = m.group(1), m.group(2), int(m.group(3))
cloudra/benchmark/gen_easy_tasks.py:276	    q = int(q_raw.replace("−", "-"))
cloudra/benchmark/gen_easy_tasks.py:277	    if ">>" in rxn:
cloudra/benchmark/gen_easy_tasks.py:278	        reactants, products = rxn.split(">>", 1)
cloudra/benchmark/gen_easy_tasks.py:279	    elif ">" in rxn:
cloudra/benchmark/gen_easy_tasks.py:280	        reactants, products = rxn.split(">", 1)
cloudra/benchmark/gen_easy_tasks.py:281	    else:
cloudra/benchmark/gen_easy_tasks.py:282	        reactants, products = rxn, ""
cloudra/benchmark/gen_easy_tasks.py:283	    r = [s for s in reactants.split(".") if s]
cloudra/benchmark/gen_easy_tasks.py:284	    p = [s for s in products.split(".") if s]
\end{lstlisting}
\clearpage
\begin{lstlisting}[firstnumber=285,numbers=left,numberstyle=\tiny\ttfamily,numbersep=5pt]
cloudra/benchmark/gen_easy_tasks.py:285	    return r, p, q, mult
cloudra/benchmark/gen_easy_tasks.py:286	
cloudra/benchmark/gen_easy_tasks.py:287	
cloudra/benchmark/gen_easy_tasks.py:288	def _parse_barrier(cell: str) -> tuple[float | None, BarrierType | None, VerificationStatus, str]:
cloudra/benchmark/gen_easy_tasks.py:289	    """Return (barrier_kcal, barrier_type, verification, doi).
cloudra/benchmark/gen_easy_tasks.py:290	
cloudra/benchmark/gen_easy_tasks.py:291	    能垒锚点只取引用括号之前的文本——括号里是文献出处（年份、卷页、DOI），
cloudra/benchmark/gen_easy_tasks.py:292	    从整格抓第一个数字会把年份（2023）或键长（2.4 Å）当成能垒。百分数
cloudra/benchmark/gen_easy_tasks.py:293	    （"能垒降 ~30%"）跳过。找不到锚点数字时返回 None（诚实缺测），并拒绝
cloudra/benchmark/gen_easy_tasks.py:294	    >200 kcal/mol 的离谱值（fail-loud 防回归）。类型只认单元格里显式写的
cloudra/benchmark/gen_easy_tasks.py:295	    Ea/ΔH‡/ΔE‡/ΔG‡ 记号；没写就是 None，不默认贴 Ea 标签。"""
cloudra/benchmark/gen_easy_tasks.py:296	    vtag = VERIFY_RE.search(cell)
cloudra/benchmark/gen_easy_tasks.py:297	    verification = _VERIFY_MAP.get(vtag.group(1), VerificationStatus.UNVERIFIED) if vtag else VerificationStatus.UNVERIFIED
cloudra/benchmark/gen_easy_tasks.py:298	    doi_m = re.search(r"DOI\s+([0-9.]+/\S*?)(?=[\s，。；）]|$)", cell)
cloudra/benchmark/gen_easy_tasks.py:299	    doi = doi_m.group(1).rstrip("）") if doi_m else ""
cloudra/benchmark/gen_easy_tasks.py:300	    # DOI 可以合法含括号（如 10.1016/s0009-2614(96)01253-5）：只剥不平衡的尾部 ")"
cloudra/benchmark/gen_easy_tasks.py:301	    while doi.endswith(")") and doi.count(")") > doi.count("("):
cloudra/benchmark/gen_easy_tasks.py:302	        doi = doi[:-1]
cloudra/benchmark/gen_easy_tasks.py:303	
cloudra/benchmark/gen_easy_tasks.py:304	    btype: BarrierType | None = None
cloudra/benchmark/gen_easy_tasks.py:305	    for k, v in _BARRIER_TYPE_MAP.items():
cloudra/benchmark/gen_easy_tasks.py:306	        if k in cell:
cloudra/benchmark/gen_easy_tasks.py:307	            btype = v
cloudra/benchmark/gen_easy_tasks.py:308	            break
cloudra/benchmark/gen_easy_tasks.py:309	
cloudra/benchmark/gen_easy_tasks.py:310	    # 引用信息从第一个括号开始；锚点在括号之前。范围取第一个数（下界）。
cloudra/benchmark/gen_easy_tasks.py:311	    head = re.split(r"[（(]", cell, maxsplit=1)[0]
cloudra/benchmark/gen_easy_tasks.py:312	    barrier = None
cloudra/benchmark/gen_easy_tasks.py:313	    for num_m in re.finditer(r"(\d+(?:\.\d+)?)(%?)", head):
cloudra/benchmark/gen_easy_tasks.py:314	        if num_m.group(2) == "%":
cloudra/benchmark/gen_easy_tasks.py:315	            continue  # "能垒降 ~30%" 是降幅不是锚点
cloudra/benchmark/gen_easy_tasks.py:316	        barrier = float(num_m.group(1))
cloudra/benchmark/gen_easy_tasks.py:317	        break
cloudra/benchmark/gen_easy_tasks.py:318	    if barrier is not None and not (0.0 <= barrier <= 200.0):
cloudra/benchmark/gen_easy_tasks.py:319	        raise ValueError(f"implausible barrier {barrier} parsed from cell: {cell!r}")
cloudra/benchmark/gen_easy_tasks.py:320	    return barrier, btype, verification, doi
cloudra/benchmark/gen_easy_tasks.py:321	
cloudra/benchmark/gen_easy_tasks.py:322	
cloudra/benchmark/gen_easy_tasks.py:323	def _parse_comfort(cell: str) -> tuple[list[str], int | None]:
cloudra/benchmark/gen_easy_tasks.py:324	    elem_m = ELEMENTS_RE.match(cell)
cloudra/benchmark/gen_easy_tasks.py:325	    elements = elem_m.group(1).split(",") if elem_m else []
cloudra/benchmark/gen_easy_tasks.py:326	    heavy_m = HEAVY_RE.search(cell)
cloudra/benchmark/gen_easy_tasks.py:327	    heavy = int(heavy_m.group(1)) if heavy_m else None
cloudra/benchmark/gen_easy_tasks.py:328	    return elements, heavy
cloudra/benchmark/gen_easy_tasks.py:329	
cloudra/benchmark/gen_easy_tasks.py:330	
cloudra/benchmark/gen_easy_tasks.py:331	def _build_graph(task_id: str, reactants: list[str], products: list[str], step_type: str) -> GoldMechanismGraph:
cloudra/benchmark/gen_easy_tasks.py:332	    """Build a minimal single-step mechanism graph (reactant -> TS -> product).
cloudra/benchmark/gen_easy_tasks.py:333	
cloudra/benchmark/gen_easy_tasks.py:334	    Easy-tier tasks are all single-step; degenerate processes (product ==
cloudra/benchmark/gen_easy_tasks.py:335	    reactant) still get distinct R/P node ids but identical SMILES.
cloudra/benchmark/gen_easy_tasks.py:336	    """
\end{lstlisting}
\clearpage
\begin{lstlisting}[firstnumber=337,numbers=left,numberstyle=\tiny\ttfamily,numbersep=5pt]
cloudra/benchmark/gen_easy_tasks.py:337	    r_node = MechanismNode(id=f"{task_id}_R", role=NodeRole.REACTANT, smiles=reactants, label="reactant")
cloudra/benchmark/gen_easy_tasks.py:338	    p_node = MechanismNode(
cloudra/benchmark/gen_easy_tasks.py:339	        id=f"{task_id}_P",
cloudra/benchmark/gen_easy_tasks.py:340	        role=NodeRole.PRODUCT,
cloudra/benchmark/gen_easy_tasks.py:341	        smiles=products if products else reactants,
cloudra/benchmark/gen_easy_tasks.py:342	        label="product",
cloudra/benchmark/gen_easy_tasks.py:343	    )
cloudra/benchmark/gen_easy_tasks.py:344	    ts_node = MechanismNode(
cloudra/benchmark/gen_easy_tasks.py:345	        id=f"{task_id}_TS", role=NodeRole.TRANSITION_STATE, smiles=[], label="TS"
cloudra/benchmark/gen_easy_tasks.py:346	    )
cloudra/benchmark/gen_easy_tasks.py:347	    # gold 边必须携带键变化（未登记 KeyError）：索引引用 query 展示的映射号，
cloudra/benchmark/gen_easy_tasks.py:348	    # 否则"agent 勤勉填了键变化"反而会因对称差被扣分（评分倒挂）
cloudra/benchmark/gen_easy_tasks.py:349	    breaks, forms = _BOND_CHANGES[task_id]
cloudra/benchmark/gen_easy_tasks.py:350	    edge = MechanismEdge(
cloudra/benchmark/gen_easy_tasks.py:351	        id=f"{task_id}_E1",
cloudra/benchmark/gen_easy_tasks.py:352	        **{"from": r_node.id, "to": p_node.id},
cloudra/benchmark/gen_easy_tasks.py:353	        via=ts_node.id,
cloudra/benchmark/gen_easy_tasks.py:354	        step_type=step_type,
cloudra/benchmark/gen_easy_tasks.py:355	        bond_changes=BondChangeSet(**{"break": breaks}, form=forms),
cloudra/benchmark/gen_easy_tasks.py:356	    )
cloudra/benchmark/gen_easy_tasks.py:357	    return GoldMechanismGraph(nodes=[r_node, ts_node, p_node], edges=[edge])
cloudra/benchmark/gen_easy_tasks.py:358	
cloudra/benchmark/gen_easy_tasks.py:359	
cloudra/benchmark/gen_easy_tasks.py:360	def _query(task_id: str, name: str, reactants: list[str], products: list[str], q: int, mult: int, conditions: str) -> str:
cloudra/benchmark/gen_easy_tasks.py:361	    rs = ".".join(reactants)
cloudra/benchmark/gen_easy_tasks.py:362	    ps = ".".join(products) if products else rs
cloudra/benchmark/gen_easy_tasks.py:363	    mapped = _MAPPED_REACTANTS[task_id]  # 原子映射反应物 SMILES；未登记 KeyError
cloudra/benchmark/gen_easy_tasks.py:364	    kind = "process" if task_id.startswith("M") else "reaction"
cloudra/benchmark/gen_easy_tasks.py:365	    # 只给任务与交付协议，不教练式地点名 TS/IRC/频率——"determine the
cloudra/benchmark/gen_easy_tasks.py:366	    # mechanism" 对合格的 agent 已隐含这些要求。证据文件的 smiles 注释是
cloudra/benchmark/gen_easy_tasks.py:367	    # 交付协议的一部分（评分器据此把图元素绑定到具体计算产物）。映射 SMILES
cloudra/benchmark/gen_easy_tasks.py:368	    # 只定义 bond_changes 的原子编号口径；节点与证据注释仍用无映射 SMILES。
cloudra/benchmark/gen_easy_tasks.py:369	    return (
cloudra/benchmark/gen_easy_tasks.py:370	        f"Determine the mechanism of the following {kind}: {name}. "
cloudra/benchmark/gen_easy_tasks.py:371	        f"Reactant SMILES: {rs}; product SMILES: {ps}; total charge {q}, multiplicity {mult}; "
cloudra/benchmark/gen_easy_tasks.py:372	        f"conditions: {conditions}. "
cloudra/benchmark/gen_easy_tasks.py:373	        f"Atom-mapped reactant SMILES: {mapped}. "
cloudra/benchmark/gen_easy_tasks.py:374	        "The 1-based map numbers in it are the indexing convention for `bond_changes` "
cloudra/benchmark/gen_easy_tasks.py:375	        "(keep using the plain SMILES above for graph nodes and evidence annotations). "
cloudra/benchmark/gen_easy_tasks.py:376	        "You must maintain `mechanism_graph.json` in the sandbox root recording the mechanism graph: "
cloudra/benchmark/gen_easy_tasks.py:377	        "nodes (reactant/product/transition_state, with SMILES) and elementary-step edges "
cloudra/benchmark/gen_easy_tasks.py:378	        "(with `via` pointing to the TS node id). Start the graph at the reactant complex "
cloudra/benchmark/gen_easy_tasks.py:379	        "(the pre-reaction complex), not at the separated reactants. The scorer checks that "
cloudra/benchmark/gen_easy_tasks.py:380	        "your graph covers the reference mechanism as a subgraph: extra nodes or edges beyond "
cloudra/benchmark/gen_easy_tasks.py:381	        "the reference are never penalized, so record every step you actually investigated. "
cloudra/benchmark/gen_easy_tasks.py:382	        "Every elementary-step edge must carry "
cloudra/benchmark/gen_easy_tasks.py:383	        "`bond_changes`: `break`/`form` lists of [i, j] atom-map index pairs, where [i, j] "
cloudra/benchmark/gen_easy_tasks.py:384	        "belongs to `break` if the bond order between atoms i and j decreases over the step "
cloudra/benchmark/gen_easy_tasks.py:385	        "(including full cleavage) and to `form` if it increases (including newly formed bonds); "
cloudra/benchmark/gen_easy_tasks.py:386	        "a purely conformational or inversion step carries two empty lists. "
cloudra/benchmark/gen_easy_tasks.py:387	        "The scorer reads this file forensically; "
cloudra/benchmark/gen_easy_tasks.py:388	        "a missing or unparseable file severely impacts the score. "
\end{lstlisting}
\clearpage
\begin{lstlisting}[firstnumber=389,numbers=left,numberstyle=\tiny\ttfamily,numbersep=5pt]
cloudra/benchmark/gen_easy_tasks.py:389	        "Every computational evidence file you produce for a species or step (QM input or log) "
cloudra/benchmark/gen_easy_tasks.py:390	        "must carry a `smiles: <SMILES>` comment line identifying the species, so the scorer "
cloudra/benchmark/gen_easy_tasks.py:391	        "can bind each graph element to concrete evidence. "
cloudra/benchmark/gen_easy_tasks.py:392	        "In your final summary you must report the computed barrier in kcal/mol for the "
cloudra/benchmark/gen_easy_tasks.py:393	        "rate-determining step, relative to the reactant complex (the pre-reaction complex, "
cloudra/benchmark/gen_easy_tasks.py:394	        "e.g. the ion-dipole complex of a bimolecular reaction; for a unimolecular process "
cloudra/benchmark/gen_easy_tasks.py:395	        "this is simply the reactant minimum). "
cloudra/benchmark/gen_easy_tasks.py:396	        "Run the computations with ORCA or Gaussian: the scorer parses only these two "
cloudra/benchmark/gen_easy_tasks.py:397	        "programs' log formats, so output from any other program cannot count as evidence. "
cloudra/benchmark/gen_easy_tasks.py:398	        "Allowed computational methods are limited to MACE-OMol25 and xTB; "
cloudra/benchmark/gen_easy_tasks.py:399	        "calls beyond the whitelist are intercepted."
cloudra/benchmark/gen_easy_tasks.py:400	    )
cloudra/benchmark/gen_easy_tasks.py:401	
cloudra/benchmark/gen_easy_tasks.py:402	
cloudra/benchmark/gen_easy_tasks.py:403	def parse_easy_md(src: Path) -> list[TaskSpec]:
cloudra/benchmark/gen_easy_tasks.py:404	    tasks: list[TaskSpec] = []
cloudra/benchmark/gen_easy_tasks.py:405	    for line in src.read_text(encoding="utf-8").splitlines():
cloudra/benchmark/gen_easy_tasks.py:406	        m = ROW_RE.match(line)
cloudra/benchmark/gen_easy_tasks.py:407	        if not m:
cloudra/benchmark/gen_easy_tasks.py:408	            continue
cloudra/benchmark/gen_easy_tasks.py:409	        cells = _split_cells(line)
cloudra/benchmark/gen_easy_tasks.py:410	        # expect: [#, name, smiles_cell, conditions, mechanism, barrier, comfort, reference]
cloudra/benchmark/gen_easy_tasks.py:411	        if len(cells) < 8:
cloudra/benchmark/gen_easy_tasks.py:412	            continue
cloudra/benchmark/gen_easy_tasks.py:413	        task_id = cells[0]
cloudra/benchmark/gen_easy_tasks.py:414	        name = _NAME_EN[task_id]  # 题名英译；未登记直接 KeyError
cloudra/benchmark/gen_easy_tasks.py:415	        smiles_cell, conditions_cell, _mechanism_cell, barrier_cell, comfort_cell, _ref_cell = (
cloudra/benchmark/gen_easy_tasks.py:416	            cells[2], cells[3], cells[4], cells[5], cells[6], cells[7],
cloudra/benchmark/gen_easy_tasks.py:417	        )
cloudra/benchmark/gen_easy_tasks.py:418	        parsed = _parse_smiles_cell(smiles_cell)
cloudra/benchmark/gen_easy_tasks.py:419	        if parsed is None:
cloudra/benchmark/gen_easy_tasks.py:420	            continue
cloudra/benchmark/gen_easy_tasks.py:421	        reactants, products, q, mult = parsed
cloudra/benchmark/gen_easy_tasks.py:422	        barrier, btype, verification, doi = _parse_barrier(barrier_cell)
cloudra/benchmark/gen_easy_tasks.py:423	        # 口径统一 CENTRAL（见 _BARRIER_REFERENCE 注释）；锚点属渐近线口径的题
cloudra/benchmark/gen_easy_tasks.py:424	        # 按 _BARRIER_FRAME_UNVERIFIED 登记降为 unverified，数值保留不编造
cloudra/benchmark/gen_easy_tasks.py:425	        barrier_ref = _BARRIER_REFERENCE.get(task_id, BarrierReference.CENTRAL) if barrier is not None else None
cloudra/benchmark/gen_easy_tasks.py:426	        if task_id in _BARRIER_FRAME_UNVERIFIED:
cloudra/benchmark/gen_easy_tasks.py:427	            verification = VerificationStatus.UNVERIFIED
cloudra/benchmark/gen_easy_tasks.py:428	        elements, heavy = _parse_comfort(comfort_cell)
cloudra/benchmark/gen_easy_tasks.py:429	        # M 组第 4 列是过程类型而非条件（类型已并入 step_type 英译）；
cloudra/benchmark/gen_easy_tasks.py:430	        # 映射表未收录的条件/概要直接 KeyError 报错，防止静默漏译。
cloudra/benchmark/gen_easy_tasks.py:431	        if task_id.startswith("M"):
cloudra/benchmark/gen_easy_tasks.py:432	            conditions = "gas phase"
cloudra/benchmark/gen_easy_tasks.py:433	        else:
cloudra/benchmark/gen_easy_tasks.py:434	            conditions = _CONDITIONS_EN[conditions_cell]
cloudra/benchmark/gen_easy_tasks.py:435	        step_type = _STEP_TYPE_EN[task_id]
cloudra/benchmark/gen_easy_tasks.py:436	
cloudra/benchmark/gen_easy_tasks.py:437	        graph = _build_graph(task_id, reactants, products, step_type)
cloudra/benchmark/gen_easy_tasks.py:438	        spec = TaskSpec(
cloudra/benchmark/gen_easy_tasks.py:439	            task_id=task_id,
cloudra/benchmark/gen_easy_tasks.py:440	            tier=Tier.EASY,
\end{lstlisting}
\clearpage
\begin{lstlisting}[firstnumber=441,numbers=left,numberstyle=\tiny\ttfamily,numbersep=5pt]
cloudra/benchmark/gen_easy_tasks.py:441	            family=task_id[0],
cloudra/benchmark/gen_easy_tasks.py:442	            name=name,
cloudra/benchmark/gen_easy_tasks.py:443	            prompt=PromptSpec(query=_query(task_id, name, reactants, products, q, mult, conditions)),
cloudra/benchmark/gen_easy_tasks.py:444	            system=SystemSpec(
cloudra/benchmark/gen_easy_tasks.py:445	                reactants=reactants,
cloudra/benchmark/gen_easy_tasks.py:446	                products=products,
cloudra/benchmark/gen_easy_tasks.py:447	                charge=q,
cloudra/benchmark/gen_easy_tasks.py:448	                multiplicity=mult,
cloudra/benchmark/gen_easy_tasks.py:449	                conditions=conditions,
cloudra/benchmark/gen_easy_tasks.py:450	                elements=elements,
cloudra/benchmark/gen_easy_tasks.py:451	                max_heavy_atoms=heavy,
cloudra/benchmark/gen_easy_tasks.py:452	                closed_shell=True,
cloudra/benchmark/gen_easy_tasks.py:453	                single_step=True,
cloudra/benchmark/gen_easy_tasks.py:454	            ),
cloudra/benchmark/gen_easy_tasks.py:455	            gold=GoldSpec(
cloudra/benchmark/gen_easy_tasks.py:456	                mechanism_graph=graph,
cloudra/benchmark/gen_easy_tasks.py:457	                reference_computation=ReferenceComputation(
cloudra/benchmark/gen_easy_tasks.py:458	                    level_of_theory=(
cloudra/benchmark/gen_easy_tasks.py:459	                        "literature anchor (see reference_doi)"
cloudra/benchmark/gen_easy_tasks.py:460	                        if doi
cloudra/benchmark/gen_easy_tasks.py:461	                        else "literature anchor (curated in easy.md; no DOI on record)"
cloudra/benchmark/gen_easy_tasks.py:462	                    ),
cloudra/benchmark/gen_easy_tasks.py:463	                    barrier_kcal_mol=barrier,
cloudra/benchmark/gen_easy_tasks.py:464	                    barrier_type=btype,
cloudra/benchmark/gen_easy_tasks.py:465	                    barrier_reference=barrier_ref,
cloudra/benchmark/gen_easy_tasks.py:466	                    verification=verification,
cloudra/benchmark/gen_easy_tasks.py:467	                    reference_doi=doi,
cloudra/benchmark/gen_easy_tasks.py:468	                    computed_by_us=False,
cloudra/benchmark/gen_easy_tasks.py:469	                ),
cloudra/benchmark/gen_easy_tasks.py:470	            ),
cloudra/benchmark/gen_easy_tasks.py:471	            rubric=RubricSpec(),
cloudra/benchmark/gen_easy_tasks.py:472	        )
cloudra/benchmark/gen_easy_tasks.py:473	        tasks.append(spec)
cloudra/benchmark/gen_easy_tasks.py:474	    return tasks
cloudra/benchmark/gen_easy_tasks.py:475	
cloudra/benchmark/gen_easy_tasks.py:476	
cloudra/benchmark/gen_easy_tasks.py:477	def main() -> None:
cloudra/benchmark/gen_easy_tasks.py:478	    ap = argparse.ArgumentParser()
cloudra/benchmark/gen_easy_tasks.py:479	    ap.add_argument("--src", type=Path, default=DEFAULT_SRC)
cloudra/benchmark/gen_easy_tasks.py:480	    ap.add_argument("--out", type=Path, default=DEFAULT_OUT)
cloudra/benchmark/gen_easy_tasks.py:481	    args = ap.parse_args()
cloudra/benchmark/gen_easy_tasks.py:482	
cloudra/benchmark/gen_easy_tasks.py:483	    tasks = parse_easy_md(args.src)
cloudra/benchmark/gen_easy_tasks.py:484	    # 先断言数量再落盘：计数回归时不会留下已写出的不一致产物
cloudra/benchmark/gen_easy_tasks.py:485	    assert len(tasks) == 30, f"expected 30 tasks, got {len(tasks)}"
cloudra/benchmark/gen_easy_tasks.py:486	    args.out.mkdir(parents=True, exist_ok=True)
cloudra/benchmark/gen_easy_tasks.py:487	    for t in tasks:
cloudra/benchmark/gen_easy_tasks.py:488	        out = args.out / f"{t.task_id}.json"
cloudra/benchmark/gen_easy_tasks.py:489	        out.write_text(
cloudra/benchmark/gen_easy_tasks.py:490	            json.dumps(t.model_dump(mode="json", by_alias=True), ensure_ascii=False, indent=2),
cloudra/benchmark/gen_easy_tasks.py:491	            encoding="utf-8",
cloudra/benchmark/gen_easy_tasks.py:492	        )
\end{lstlisting}
\clearpage
\end{document}
